\sloppy % Membantu LaTeX memecah kata dengan lebih fleksibel
\onehalfspacing % Terapkan spasi 1,5 seperti MS Word

\setcounter{page}{78}

% ============================================================
% BAB IV: PERANCANGAN DAN IMPLEMENTASI SISTEM
% ============================================================

\vspace{2cm}
\begin{center}
{\fontsize{14}{16.8}\selectfont\textbf{BAB IV. Perancangan dan Implementasi Sistem}}\\[1em]
\end{center}
\addcontentsline{toc}{section}{BAB IV PERANCANGAN DAN IMPLEMENTASI SISTEM}
\vspace{2em}

% Pengaturan format penomoran untuk Chapter IV
\renewcommand{\thesection}{\Roman{section}}
\renewcommand{\thesubsection}{\thesection.\arabic{subsection}}
\renewcommand{\thesubsubsection}{\thesubsection.\arabic{subsubsection}}


% Pengaturan format penomoran tabel mengikuti section
\counterwithin{table}{section}
\renewcommand{\thetable}{\thesection.\arabic{table}}

% Pengaturan format penomoran gambar mengikuti section
\counterwithin{figure}{section}
\renewcommand{\thefigure}{\thesection.\arabic{figure}}

% Pengaturan format penomoran persamaan mengikuti section
\counterwithin{equation}{section}
\renewcommand{\theequation}{\thesection.\arabic{equation}}

% Section utama BAB IV
\setcounter{section}{4}
% \section{Perancangan dan Implementasi Sistem}
\label{sec:perancangan-implementasi}

Bab ini menguraikan rancangan solusi teknis yang diusulkan untuk restorasi dokumen terdegradasi dengan pendekatan \textit{Generative Adversarial Network} yang dimodifikasi secara khusus untuk meningkatkan keterbacaan teks oleh model pengenalan tulisan tangan. Berdasarkan analisis kebutuhan yang telah ditetapkan pada Bab III (8 kebutuhan fungsional dan 4 kebutuhan non-fungsional), bab ini menyajikan perancangan dan implementasi sistem secara komprehensif mulai dari arsitektur teoretis hingga protokol eksperimental yang dapat direproduksi.

Bab ini disusun dalam tujuh bagian utama yang mengikuti alur logis dari desain hingga evaluasi. Bagian pertama (\S\ref{subsec:rancangan-arsitektur}) menguraikan rancangan arsitektur framework GAN-HTR, mencakup justifikasi pemilihan komponen, spesifikasi detail Generator U-Net Enhanced (21.8M parameter), Frozen Recognizer berbasis Transformer (27.86M parameter), Diskriminator Dual-Modal dengan fusi cross-modal bilateral (17.4M parameter), dan formulasi fungsi loss multi-komponen yang dioptimalkan. Bagian kedua (\S\ref{subsec:lingkungan-eksperimen}) menjelaskan infrastruktur komputasi dan lingkungan eksperimen yang digunakan untuk implementasi, termasuk spesifikasi perangkat keras, konfigurasi software stack, dan mekanisme reproduktibilitas. Bagian ketiga (\S\ref{subsec:pengaturan-eksperimental}) menguraikan pengaturan eksperimental komprehensif, meliputi konstruksi dataset semi-sintetis realistis dan dataset dokumen historis nyata ANRI untuk evaluasi kuantitatif dan kualitatif. Bagian keempat (\S\ref{subsec:implementasi-pipeline}) menjelaskan implementasi pipeline persiapan data untuk training recognizer dan degradasi sintetis untuk training GAN. Bagian kelima (\S\ref{subsec:detail-implementasi-training}) merinci strategi dan protokol training, termasuk pendekatan curriculum learning tiga fase, prosedur adversarial bergantian, konfigurasi optimisasi, dan teknik regularisasi. Bagian keenam (\S\ref{subsec:metode-baseline-strategi}) memaparkan metode baseline klasik dan strategi evaluasi komparatif untuk memvalidasi keunggulan metode yang diusulkan. Bagian ketujuh (\S\ref{subsec:konfigurasi-metrik-evaluasi}) merinci konfigurasi metrik evaluasi multi-objektif (PSNR, SSIM, CER, WER) beserta spesifikasi implementasi teknis dan protokol eksperimen ablasi untuk validasi empiris pada Bab V.

\subsection{Rancangan Arsitektur}
\label{subsec:rancangan-arsitektur}

Rancangan solusi bertujuan membangun model generatif yang dioptimalkan untuk restorasi visual dan keterbacaan tekstual simultan. Tiga inovasi kunci: (1) frozen recognizer memberikan gradien keterbacaan stabil tanpa catastrophic forgetting. (2) Fungsi loss multi-komponen menyeimbangkan kualitas visual (L1 + perceptual + adversarial) dan keterbacaan teks (CTC loss). (3) Diskriminator dual-modal (CNN+BiLSTM) mengevaluasi koherensi visual-tekstual. Ketiga inovasi ini divalidasi empiris pada Bab V, dengan pemetaan kebutuhan sistem di Bab III Section 3.2.3.

\subsubsection{Pemilihan Arsitektur Baseline dan Hipotesis Inovasi}
\label{subsubsec:justifikasi-arsitektur}

Pemilihan arsitektur framework GAN-HTR yang diusulkan didasarkan pada analisis kesenjangan dari metode \textit{state-of-the-art} yang telah dikaji pada Bab II (Tinjauan Pustaka). Bagian ini menguraikan: (1) justifikasi komponen baseline berdasarkan analisis kesenjangan literatur, dan (2) hipotesis inovasi yang akan divalidasi melalui studi ablasi empiris pada Bab V. Tabel \ref{tab:arsitektur-komparasi} membandingkan karakteristik arsitektur metode yang ada dengan framework yang diusulkan.

\begin{table}[H]
\caption[Komparasi arsitektur framework dengan metode state-of-the-art]{Komparasi arsitektur framework yang diusulkan dengan metode state-of-the-art. Perbandingan sistematis berdasarkan komponen generator, diskriminator, strategi integrasi HTR, dan keterbatasan masing-masing metode. Framework GAN-HTR yang diusulkan mengintegrasikan diskriminator dual-modal dengan frozen recognizer dan CTC loss untuk optimasi simultan kualitas visual dan keterbacaan teks.}
\label{tab:arsitektur-komparasi}
\footnotesize
\centering
\begin{tabularx}{\textwidth}{|X|X|X|X|X|}
\hline
\textbf{Metode} & Generator & \textbf{Diskriminator} & Integrasi HTR & Keterbatasan \\
\hline
\multicolumn{5}{|l|}{Metode Existing (Bab II)} \\
\hline
DE-GAN & U-Net & CNN (visual-only) & Post-training eval & Single-modal, no HTR loss \\
\hline
DocEnTr & Transformer & - (no adversarial) & Post-training eval & Over-smoothing, no texture \\
\hline
Text-DIAE & CNN-based & - (no adversarial) & Post-training eval & No pelatihan adversarial \\
\hline
ERB-MultiTask & U-Net & CNN (visual-only) & Joint training & Unstable, catastrophic forgetting \\
\hline
\multicolumn{5}{|l|}{Framework yang Diusulkan} \\
\hline
GAN-HTR & U-Net Enhanced & Dual-Modal & Frozen recognizer & Kontribusi dual-modal \\
(penelitian ini) &  & (CNN+LSTM) & + CTC loss & marginal ($p>0.05$) \\
\hline
\end{tabularx}
\end{table}


\subsubsection{Komponen Utama dan Interaksi Arsitektur}
\label{subsubsec:komponen-arsitektur}

Untuk mengatasi keterbatasan metode restorasi konvensional yang mengabaikan keterbacaan teks mesin, penelitian ini mengusulkan kerangka kerja GAN berorientasi HTR yang secara eksplisit mengintegrasikan pengenal teks dalam loop pelatihan. Kerangka kerja yang diusulkan, yang diilustrasikan pada Gambar~\ref{fig:arch_overview}, beroperasi pada citra baris dokumen tulisan tangan yang terdegradasi dan menghasilkan keluaran yang bersih dan dapat dibaca yang dioptimalkan untuk kualitas visual dan kinerja HTR.

Arsitektur yang diusulkan terdiri dari tiga komponen utama yang berinteraksi dalam sebuah kerangka pelatihan adversarial. Alur proses restorasi berlangsung sebagai berikut: Citra terdegradasi $I_{\text{deg}}$ dimasukkan ke Generator $G$ yang menghasilkan citra terestorasi $I_{\text{gen}}$. Diskriminator dual-modal $D$ menerima pasangan $(I_{\text{gt}}, \text{teks}_{\text{gt}})$ sebagai sampel nyata dan $(I_{\text{gen}}, \text{teks}_{\text{pred}})$ sebagai sampel palsu untuk pembelajaran adversarial. Secara paralel, pengenal HTR beku $R$ mengekstrak fitur dari $I_{\text{gen}}$ untuk memastikan keterbacaan teks melalui \textit{CTC loss} dan \textit{recognition feature loss}.

Tiga komponen utama kerangka kerja:
\begin{enumerate}[nosep]
    \item Generator (G): Arsitektur U-Net Enhanced yang memetakan citra terdegradasi $I_{\text{deg}}$ ke citra yang direstorasi $I_{\text{gen}}$. Generator dilengkapi \textit{Residual Dense Blocks} (RDB) dan \textit{Attention Gates} untuk preservasi detail teks halus sambil menghilangkan noise latar belakang. Memenuhi FR-1 (Restorasi Gambar) dan FR-2 (Penanganan Degradasi).
    
    \item Frozen Recognizer (R): Recognizer berbasis Transformer praterlatih dengan bobot tetap yang menyediakan gradien sadar HTR tanpa risiko catastrophic forgetting. Recognizer mencapai CER 33.72\% pada dataset validasi awal ANRI ($n$=948), dan CER 26.57\% (set validasi, $n$=710) serta 34.1\% (set uji, $n$=712) pada citra bersih \textit{ground truth} penelitian ini---perbedaan mencerminkan variasi kompleksitas teks antar split dataset. Strategi pembekuan mengatasi joint training instability dengan memisahkan optimasi visual dari evaluasi tekstual. Memenuhi FR-3 (Integrasi HTR) dan NFR-2 (Keterbacaan Teks).
    
    \item Diskriminator Dual-Modal (D): Diskriminator dengan jalur CNN dan BiLSTM paralel yang menilai kualitas visual dan koherensi teks secara bersamaan. Mekanisme \textit{bilateral cross-attention fusion} memungkinkan validasi konsistensi visual-tekstual untuk menghindari artefak yang merusak struktur karakter. Mendukung NFR-1 (Kualitas Visual).
\end{enumerate}

\begin{figure}[htbp!]
\centering
\includegraphics[width=1.35\textwidth, height=0.85\textheight, keepaspectratio]{/home/lambda_one/tesis/GAN-HTR-ORI/docRestoration/Paper/drawio/main_framework.png}
\caption[Kerangka kerja GAN-HTR dengan diskriminator dual-modal]{Kerangka kerja GAN-HTR dengan diskriminator dual-modal dan \textit{frozen} HTR recognizer.\\
Komponen utama: (1) Generator U-Net Enhanced, (2) Diskriminator Dual-Modal (CNN+BiLSTM),\\
(3) Frozen HTR Recognizer untuk CTC loss.}
\label{fig:arch_overview}
\end{figure}

Ketiga komponen berinteraksi dalam siklus pelatihan adversarial: Generator menghasilkan gambar restorasi dari input terdegradasi, frozen recognizer mengevaluasi keterbacaan melalui prediksi karakter dan CTC loss, sementara diskriminator menilai koherensi visual-tekstual dari pasangan gambar dan representasi teks. Proses pelatihan mengadopsi strategi adversarial bergantian dengan curriculum learning tiga fase (detail pada Subsection~\ref{subsec:detail-implementasi-training}): (1) Pembaruan diskriminator---$D$ dilatih untuk membedakan pasangan $(I_{\text{gt}}, \text{teks}_{\text{gt}})$ nyata dari pasangan $(I_{\text{gen}}, \text{teks}_{\text{pred}})$ yang dihasilkan, dengan supervisi dual-modal pada domain visual dan tekstual secara bersamaan. (2) Pembaruan Generator---$G$ dioptimalkan untuk menipu $D$ sambil meminimalkan $\mathcal{L}_{\text{pixel}}$ terhadap \textit{ground truth}, $\mathcal{L}_{\text{perc}}$ untuk preservasi struktur perseptual, $\mathcal{L}_{\text{rec-feat}}$ untuk konsistensi fitur HTR, dan $\mathcal{L}_{\text{CTC}}$ untuk keterbacaan teks eksplisit (konfigurasi optimal pada Persamaan~\ref{eq:IV.13}).\sloppy

\subsubsection{Arsitektur Generator: U-Net Enhanced dengan Blok Residual}
\label{subsubsec:arsitektur-generator}

Penelitian ini mengadopsi arsitektur U-Net Enhanced yang menggabungkan empat teknik restorasi dokumen berkualitas tinggi. Berbeda dari U-Net konvensional yang hanya menggunakan blok konvolusional standar, arsitektur Enhanced mengintegrasikan: (1) \textit{Residual Dense Blocks} (RDB)~\cite{zhang2018residual} untuk ekstraksi fitur hierarkis dengan aliran gradien yang lebih baik, (2) \textit{Convolutional Block Attention Module} (CBAM)~\cite{woo2018cbam} untuk rekalibrasi fitur adaptif pada dimensi kanal dan spasial, (3) \textit{Multi-Scale Feature Pyramid} (MSFP)~\cite{lin2017fpn} dengan konvolusi terdilasi~\cite{chen2017deeplab} pada bottleneck untuk agregasi konteks multi-resolusi, dan (4) \textit{Attention Gates}~\cite{oktay2018attention} pada decoder untuk fokus selektif pada wilayah pembawa teks sambil menekan amplifikasi noise latar belakang. Kombinasi teknik-teknik ini memungkinkan pelestarian goresan tipis dan tanda diakritik paleografi sambil menghilangkan degradasi kompleks seperti bleed-through, noda, dan derau.

\paragraph{Spesifikasi Arsitektur Generator}

Tabel~\ref{tab:generator_arch} merinci arsitektur lengkap generator. Encoder secara progresif melakukan downsampling pada citra masukan melalui lima tahap, mengekstrak fitur multi-skala pada resolusi dari $1024 \times 128$ hingga $32 \times 4$ (format: lebar $\times$ tinggi). Setiap blok encoder terdiri dari:

\begin{itemize}
    \item Residual Dense Block (RDB): Tiga lapisan konvolusional densely-connected (kernel $3 \times 3$) dengan growth rate 32, menghasilkan fitur kaya melalui fusi lokal dari semua lapisan sebelumnya. Berbeda dari ResNet~\cite{he2016deep} standar yang hanya memiliki satu koneksi shortcut, RDB menggunakan koneksi padat yang memungkinkan setiap lapisan mengakses fitur dari semua lapisan sebelumnya, meningkatkan kapasitas representasi dan aliran gradien.
    \item Convolutional Block Attention Module (CBAM): Mekanisme atensi sekuensial yang terdiri dari: (a) atensi kanal dengan rasio reduksi 8, menggunakan shared MLP untuk jalur global average pooling dan global max pooling, dan (b) atensi spasial dengan kernel $7 \times 7$, menggabungkan channel-wise average dan max pooling untuk menghasilkan peta atensi fokus. CBAM memungkinkan jaringan secara adaptif memfokuskan pada kanal fitur yang informatif dan wilayah spasial yang relevan.
    \item Batch normalization (momentum~$=0.9$, $\epsilon=10^{-5}$) dan aktivasi LeakyReLU ($\alpha=0.2$)
    \item Koneksi residual untuk propagasi gradien efektif pada arsitektur dalam
\end{itemize}

\textit{Bottleneck} pada resolusi $4 \times 32$ (tinggi $\times$ lebar) menggunakan \textit{Multi-Scale Feature Pyramid} (MSFP)~\cite{lin2017fpn} dengan tiga skala konvolusi terdilasi (dilation rates 1, 2, 4)~\cite{chen2017deeplab} untuk menangkap konteks multi-resolusi, diikuti RDB 512 kanal dan CBAM untuk penyempurnaan fitur global. \textit{Decoder} mencerminkan encoder dengan lima tahap upsampling, masing-masing dilengkapi gerbang atensi untuk fokus pada wilayah teks:

\begin{equation}
\alpha = \sigma(W_g^T g + W_x^T x + b)
\label{eq:IV.1}
\end{equation}

\noindent di mana $g \in \mathbb{R}^{H \times W \times C}$ adalah fitur decoder dari skala yang lebih kasar (sinyal gating), $x \in \mathbb{R}^{H \times W \times C}$ adalah fitur koneksi langsung dari encoder, $W_g$ dan $W_x$ adalah matriks proyeksi yang dapat dipelajari, $b$ adalah bias, dan $\sigma$ adalah fungsi aktivasi sigmoid. Koefisien atensi $\alpha \in [0,1]^{H \times W \times C}$ kemudian dikalikan elemen-demi-elemen dengan koneksi langsung $x$ untuk menekan fitur latar belakang yang tidak relevan sambil memperkuat wilayah pembawa teks: $\hat{x} = \alpha \odot x$. Mekanisme ini memungkinkan decoder secara adaptif memilih fitur resolusi tinggi yang relevan dari encoder.

% Tabel Arsitektur Generator
\begin{table}[H]
\caption[Arsitektur Generator U-Net Enhanced dengan spesifikasi detail]{Arsitektur Generator U-Net Enhanced. Spesifikasi detail lima tahap encoder (progressive downsampling dari $128 \times 1024$ ke $4 \times 32$), bottleneck MSFP dengan konvolusi terdilasi, dan lima tahap decoder dengan attention gates. Setiap tahap mengintegrasikan RDB+CBAM untuk ekstraksi fitur hierarkis dan atensi adaptif. Total parameter: 21.8M.}
\label{tab:generator_arch}
\small
\centering
\resizebox{\textwidth}{!}{%
\begin{tabular}{|l|l|p{3cm}|p{3cm}|r|}
\hline
\textbf{Tahap} & \textbf{Lapisan} & \textbf{Bentuk Keluaran} & \textbf{Kernel/Stride} & \textbf{Parameter} \\
\hline
Masukan & Masukan & $128 \times 1024 \times 1$ & - & - \\
\hline
\multirow{2}{*}{Encoder-0} 
& ResBlock-1 & $128 \times 1024 \times 64$ & $3 \times 3 / 1$ & 4.7K \\
& MaxPool-1 & $64 \times 512 \times 64$ & $2 \times 2 / 2$ & - \\
\hline
\multirow{2}{*}{Encoder-1}
& ResBlock-2 & $64 \times 512 \times 128$ & $3 \times 3 / 1$ & 147.6K \\
& MaxPool-2 & $32 \times 256 \times 128$ & $2 \times 2 / 2$ & - \\
\hline
\multirow{2}{*}{Encoder-2}
& ResBlock-3 & $32 \times 256 \times 256$ & $3 \times 3 / 1$ & 590.1K \\
& MaxPool-3 & $16 \times 128 \times 256$ & $2 \times 2 / 2$ & - \\
\hline
\multirow{2}{*}{Encoder-3}
& ResBlock-4 & $16 \times 128 \times 512$ & $3 \times 3 / 1$ & 2.4M \\
& MaxPool-4 & $8 \times 64 \times 512$ & $2 \times 2 / 2$ & - \\
\hline
\multirow{2}{*}{Encoder-4}
& ResBlock-5 & $8 \times 64 \times 512$ & $3 \times 3 / 1$ & 2.4M \\
& MaxPool-5 & $4 \times 32 \times 512$ & $2 \times 2 / 2$ & - \\
\hline
Bottleneck & MSFP + RDB + CBAM & $4 \times 32 \times 512$ & - & 5.1M \\
\hline
\multirow{3}{*}{Decoder-4}
& UpSample-5 & $8 \times 64 \times 512$ & - & - \\
& AttentionGate-5 & $8 \times 64 \times 512$ & - & 262K \\
& RDB + CBAM & $8 \times 64 \times 512$ & $3 \times 3 / 1$ & 3.5M \\
\hline
\multirow{3}{*}{Decoder-3}
& UpSample-4 & $16 \times 128 \times 512$ & - & - \\
& AttentionGate-4 & $16 \times 128 \times 512$ & - & 262K \\
& RDB + CBAM & $16 \times 128 \times 512$ & $3 \times 3 / 1$ & 3.5M \\
\hline
\multirow{3}{*}{Decoder-2}
& UpSample-3 & $32 \times 256 \times 256$ & - & - \\
& AttentionGate-3 & $32 \times 256 \times 256$ & - & 66K \\
& RDB + CBAM & $32 \times 256 \times 256$ & $3 \times 3 / 1$ & 885K \\
\hline
\multirow{3}{*}{Decoder-1}
& UpSample-2 & $64 \times 512 \times 128$ & - & - \\
& AttentionGate-2 & $64 \times 512 \times 128$ & - & 16.5K \\
& RDB + CBAM & $64 \times 512 \times 128$ & $3 \times 3 / 1$ & 221K \\
\hline
\multirow{3}{*}{Decoder-0}
& UpSample-1 & $128 \times 1024 \times 64$ & - & - \\
& AttentionGate-1 & $128 \times 1024 \times 64$ & - & 4.1K \\
& RDB + CBAM & $128 \times 1024 \times 64$ & $3 \times 3 / 1$ & 55K \\
\hline
Keluaran & Conv2D + Tanh & $128 \times 1024 \times 1$ & $1 \times 1 / 1$ & 65 \\
\hline
\multicolumn{4}{|r|}{\textbf{Total Parameter:}} & \textbf{21.8M} \\
\hline
\end{tabular}
}
\end{table}

% Diagram Arsitektur Generator U-Net Enhanced
\begin{figure}[H]
\centering
\includegraphics[width=0.5\textwidth,keepaspectratio]{/home/lambda_one/tesis/GAN-HTR-ORI/docRestoration/Paper/drawio/arsitektur_generator.png}
\caption[Visualisasi arsitektur Generator U-Net Enhanced]{Arsitektur Generator U-Net Enhanced. Visualisasi struktur simetris encoder-decoder dengan lima level hierarkis. Encoder mengintegrasikan RDB+CBAM pada setiap tahap untuk ekstraksi fitur adaptif, bottleneck MSFP (512 kanal) menangkap konteks multi-skala melalui konvolusi terdilasi, decoder dengan attention gates merekonstruksi citra secara bertahap. Skip connections dengan gating mechanism mempertahankan detail spasial sambil menekan noise.}
\label{fig:unet_architecture}
\end{figure}

Statistik Arsitektur:
\begin{itemize}
    \item Total parameter: 21.8M (semua terlatih)
    \item Memory footprint: $\sim$87~MB (float32)
    \item Kedalaman: 5 tahap encoder + bottleneck + 5 tahap decoder
    \item Masukan/Keluaran: $1024 \times 128 \times 1$ (L$\times$T$\times$C)
    \item Aktivasi keluaran: \textit{Tanh} (rentang [-1, 1])
\end{itemize}

Arsitektur generator yang telah dirancang ini secara khusus dioptimalkan untuk karakteristik unik dokumen paleografi historis abad 16--18, di mana pelestarian detail goresan tipis (1--3 piksel) dan tanda diakritik halus menjadi penting untuk keterbacaan HTR, sementara degradasi kompleks seperti bleed-through dan noda organik dengan variasi spektral luas memerlukan kapasitas representasi hierarkis yang tinggi. Kombinasi RDB+CBAM pada setiap \textit{level encoder-decoder} meningkatkan kapasitas ekstraksi fitur sebesar $\sim$40\% tanpa degradasi gradien, memungkinkan model membedakan antara goresan teks autentik yang harus dilestarikan dan artefak degradasi yang harus dihilangkan---kemampuan yang tidak dimiliki arsitektur U-Net konvensional.

Gambar~\ref{fig:unet_architecture} memvisualisasikan arsitektur simetris U-Net Enhanced, di mana setiap level encoder memiliki pasangan decoder yang berkomunikasi melalui skip connections dengan attention gates. Berbeda dari U-Net konvensional~\cite{ronneberger2015unet}, integrasi RDB+CBAM pada setiap level meningkatkan kapasitas representasi sebesar $\sim$40\% tanpa degradasi gradien, sementara bottleneck MSFP menangkap konteks global multi-skala sebelum proses rekonstruksi bertahap pada decoder.

\paragraph{Rasional Desain dan Keunggulan Teknis}

Residual Dense Blocks memberikan keunggulan ganda: (1) propagasi gradien yang lebih baik melalui koneksi padat multi-jalur, memungkinkan pelatihan arsitektur dalam tanpa degradasi gradien, dan (2) penggunaan kembali fitur hierarkis di mana setiap lapisan dapat mengakses fitur dari semua lapisan sebelumnya, meningkatkan kapasitas representasi tanpa meningkatkan parameter secara substansial. CBAM memberikan selektivitas adaptif pada dimensi kanal dan spasial, memungkinkan jaringan secara otomatis memfokuskan pada kanal fitur yang informatif ("what") dan wilayah spasial yang relevan ("where"), efektif untuk menekan noise latar belakang sambil memperkuat fitur goresan teks. Multi-Scale Feature Pyramid~\cite{lin2017fpn} pada bottleneck menangkap konteks multi-resolusi melalui konvolusi terdilasi~\cite{chen2017deeplab} dengan skala berbeda, penting untuk menangani degradasi dengan ukuran bervariasi (noise titik vs noda besar). Attention Gates~\cite{oktay2018attention} memberikan selektivitas spasial pada koneksi langsung, memungkinkan decoder fokus pada wilayah pembawa teks sambil menekan amplifikasi derau---penting untuk melestarikan goresan tipis dan tanda diakritik paleografi. Kombinasi teknik-teknik ini menghasilkan arsitektur yang mampu mempertahankan detail halus sambil menghilangkan degradasi kompleks.

\subsubsection{Integrasi Frozen Recognizer HTR}
\label{subsubsec:frozen-recognizer}

Penelitian ini mengintegrasikan recognizer teks praterlatih dengan strategi pembekuan bobot (frozen recognizer) untuk memberikan gradien sadar-teks yang stabil tanpa risiko catastrophic forgetting. Berbeda dengan pendekatan joint training pada ERB-MultiTask~\cite{erb2021}, komponen Recognizer (R) berfungsi sebagai evaluator objektif yang menghasilkan distribusi probabilitas karakter untuk: (1) input teks bagi Diskriminator dual-modal dalam menilai koherensi visual-tekstual, dan (2) kalkulasi \textit{CTC Loss} dan \textit{Recognition Feature Loss} untuk mengukur keterbacaan hasil restorasi.

\paragraph{Mekanisme Pembekuan dan Pencegahan Kompetisi Gradien}

Strategi pembekuan bobot recognizer dirancang untuk mengatasi masalah fundamental pada joint training: kompetisi gradien yang destruktif antara dua tujuan yang berlawanan. Dalam joint training, Generator berusaha menghasilkan citra yang mudah dibaca (untuk menurunkan CTC loss), sementara recognizer secara simultan dilatih untuk mempertahankan akurasi pada citra bersih \textit{ground truth}. Konflik ini menyebabkan gradien yang berlawanan arah mengalir ke kedua komponen, mengakibatkan ketidakstabilan pelatihan dan catastrophic forgetting pada recognizer.

Dengan membekukan bobot recognizer (total 27.86M parameter), konflik gradien ini dieliminasi sepenuhnya. Generator menerima sinyal CTC loss yang konsisten dari recognizer yang stabil, memungkinkan optimasi terfokus pada peningkatan kualitas visual citra untuk keterbacaan—bukan adaptasi recognizer pada citra berkualitas rendah. Mekanisme ini memastikan bahwa: (1) recognizer mempertahankan kemampuan pengenalan yang telah dipelajari selama pra-pelatihan, (2) generator menerima gradien sadar-teks yang objektif dan konsisten, dan (3) pelatihan mencapai stabilitas tanpa osilasi atau mode collapse.

\paragraph{Analisis Komparatif Strategi Integrasi Pengenal HTR}

Pemilihan strategi integrasi pengenal HTR merupakan keputusan arsitektural kritis yang memengaruhi stabilitas pelatihan, efisiensi komputasi, dan kualitas akhir sistem. Validasi empiris studi ablasi (Bab~V.6.3) menunjukkan bahwa frozen recognizer mempertahankan CER 31.63\% (degradasi moderat +5.06\% dari baseline pra-pelatihan 26.57\%), sementara joint training mengalami catastrophic forgetting total dengan CER 100.00\% (degradasi +73.43\%), membuktikan keefektifan strategi pembekuan untuk domain dokumen paleografi historis.

\paragraph{Justifikasi Pemilihan Strategi Pengenal Beku}

Keputusan mengadopsi strategi pengenal beku didasarkan pada empat pertimbangan komprehensif:

\begin{table}[H]
\caption[Analisis komparatif strategi integrasi pengenal HTR]{Analisis komparatif strategi integrasi pengenal HTR.\\
Evaluasi empat pendekatan (pelatihan bersama, frozen, dua tahap, baseline tanpa recognizer)\\
berdasarkan keuntungan dan keterbatasan masing-masing. Strategi frozen yang dipilih dalam penelitian ini\\
menawarkan keseimbangan optimal antara efisiensi komputasional (48 GPU-jam untuk 50 epoch)\\
dan stabilitas pelatihan tanpa risiko catastrophic forgetting.}
\label{table:recognizer_strategy_comparison}
\small
\renewcommand{\arraystretch}{1.2}
\centering
\begin{tabular}{|p{2.5cm}|p{5.2cm}|p{5.8cm}|}
\hline
\textbf{Strategi} & \textbf{Keuntungan} & \textbf{Limitasi} \\
\hline
Pelatihan bersama & Optimasi \textit{end-to-end}, adaptasi spesifik-tugas & Biaya komputasi sangat tinggi, ketidakstabilan pelatihan dari konflik gradien, risiko pelupaan katastrofik~\cite{kirkpatrick2017overcoming} \\
\hline
Beku$^\star$ & Konvergensi stabil (48 GPU-jam untuk 50 epoch), efisiensi komputasional, melestarikan kinerja pengenal (CER 33.72\%) & Adaptasi spesifik-tugas pengenal terbatas \\
\hline
Dua tahap & Optimasi seimbang, stabilitas sekuensial & Waktu pelatihan berlipat ganda, penjadwalan kompleks, memerlukan infrastruktur pelatihan pengenal \\
\hline
\textit{Baseline} (Tanpa R) & Arsitektur sederhana, komputasi minimal & Tidak ada optimasi sadar-HTR, CER sangat tinggi pada masukan terdegradasi (83.4\%) \\
\hline
\multicolumn{3}{|l|}{\footnotesize $^\star$Strategi yang dipilih dalam penelitian ini} \\
\hline
\end{tabular}
\end{table}

\begin{enumerate}[leftmargin=1.5em, labelsep=0.5em, itemindent=0pt]
\item Efisiensi komputasional \\[0.1em]
Pengenal pra-latih (27.86M parameter) tidak memerlukan komputasi gradien, sehingga hanya Generator (89.4M parameter) dan diskriminator (17.4M parameter) yang dioptimasi. Pendekatan ini mengurangi kompleksitas optimasi dan mencegah konflik gradien antar komponen. Waktu pelatihan aktual yang tercatat adalah 17.82 GPU-jam untuk 50 epoch dengan ukuran batch 2 pada GPU RTX A4000 16GB, lebih efisien dibandingkan strategi joint training yang memerlukan optimasi simultan seluruh komponen (estimasi 21 GPU-jam pada full dataset berdasarkan eksperimen ablation 20 epoch dengan subset data 2.4\% yang berjalan 30 menit).

\item Stabilitas pelatihan \\[0.1em]
Pengenal beku mencegah gradien yang bertentangan antara tujuan rekonstruksi ($\mathcal{L}_{\text{pixel}}$, $\mathcal{L}_{\text{adv}}$) dan tujuan pengenalan ($\mathcal{L}_{\text{CTC}}$, $\mathcal{L}_{\text{rec-feat}}$). Literatur menunjukkan pelatihan bersama pada GAN multi-tugas rentan terhadap kolaps modus dan konvergensi berosilasi~\cite{souibgui2022,kang2021pay}. Dengan membekukan pengenal, konflik gradien ini dieliminasi sepenuhnya, memastikan trajektori pelatihan yang stabil dan konvergen (validasi empiris pada Gambar~\ref{fig:convergence_analysis} menunjukkan tidak ada kejadian divergensi selama 50 epoch).

\item Pelestarian kinerja \\[0.1em]
Pengenal hibrida CNN-Transformer (6-\textit{layer transformer encoder} dengan \textit{CTC decoder}) yang dilatih pada dataset ANRI mencapai CER 33.72\% saat dievaluasi pada subset ANRI terpisah ($n$=948) untuk validasi awal kemampuan generalisasi pada domain paleografi Belanda historis abad 16--18. Membekukan pengenal mempertahankan kapabilitas generalisasi ini dan mencegah pelupaan katastrofik~\cite{kirkpatrick2017overcoming} yang dapat terjadi jika pengenal dilatih ulang pada pergeseran distribusi dari citra yang ditingkatkan. Kerangka kerja yang diusulkan mencapai CER 27.11\% pada set validasi (Bab~V), menunjukkan bahwa restorasi berbasis GAN meningkatkan keterbacaan melampaui kemampuan intrinsik pengenal.

\item Diferensiasi dari literatur \\[0.1em]
Berbeda dari \citet{souibgui2022} yang menggunakan pelatihan bersama pengenal-Generator (berpotensi ketidakstabilan gradien), penelitian ini membekukan pengenal praterlatih untuk stabilitas optimasi. Strategi ini mencegah konflik gradien yang terdokumentasi pada GAN multi-tugas~\cite{souibgui2022,kang2021pay}, memungkinkan Generator fokus pada restorasi visual tanpa mendegradasi pengenal. Pendekatan ini memungkinkan restorasi mencapai keterbacaan mendekati batas atas teoretis (\textit{gap} CER hanya 0.8 poin persentase dari ground truth bersih pada set uji, lihat Bab~V).
\end{enumerate}

Validasi empiris studi ablasi (Bab~V.6.3) menunjukkan bahwa frozen recognizer mempertahankan CER 31.63\% (degradasi moderat +5.06\% dari baseline pra-pelatihan 26.57\%), sementara joint training mengalami catastrophic forgetting total dengan CER 100.00\% (degradasi +73.43\%), membuktikan keefektifan strategi pembekuan untuk domain dokumen paleografi historis.

\subsubsection{Ekstraksi Fitur dan Perhitungan Loss}
\label{subsubsec:ekstraksi-fitur-recognizer}
Penelitian ini mengekstrak fitur pengenalan perantara dari recognizer beku:

\begin{equation}
\mathcal{F}_{\text{rec}}(I) = R_{\text{encoder}}(I)
\label{eq:IV.2}
\end{equation}

\noindent di mana $R_{\text{encoder}}$ mewakili bagian encoder yang dibekukan dari pengenal. Secara spesifik, digunakan keluaran dari lapisan proyeksi sebelum \textit{encoder transformer}, menghasilkan peta fitur dengan dimensi (batch, 128, 512). Lapisan ini dipilih karena mengandung fitur tingkat urutan sebelum pengodean kontekstual, memberikan gradien yang stabil untuk Generator. Fitur-fitur ini menangkap karakteristik yang relevan dengan teks seperti pola goresan, bentuk karakter, dan pengaturan spasial.

\textit{Loss} fitur pengenalan antara citra yang dihasilkan dan citra ground truth dihitung menggunakan \textit{Mean Squared Error} (MSE):

\begin{equation}
\begin{split}
\mathcal{L}_{\text{rec-feat}} &= \text{MSE}(\mathcal{F}_{\text{rec}}(I_{\text{gen}}), \mathcal{F}_{\text{rec}}(I_{\text{gt}})) \\
&= \|\mathcal{F}_{\text{rec}}(I_{\text{gen}}) - \mathcal{F}_{\text{rec}}(I_{\text{gt}})\|_2^2
\end{split}
\label{eq:IV.3}
\end{equation}

MSE dipilih karena lebih sensitif terhadap deviasi fitur besar, memberikan gradien yang lebih kuat untuk perbaikan keselarasan fitur. Ini mendorong Generator untuk menghasilkan citra yang fitur teksnya cocok dengan fitur teks dari \textit{ground truth} yang bersih, secara implisit meningkatkan keterbacaan HTR. Arsitektur recognizer yang digunakan merupakan hibrida CNN-Transformer yang terdiri dari: CNN \textit{Backbone} 8-lapisan (4 tahap, 64$\rightarrow$512 kanal), \textit{Transformer Encoder} 6-lapisan (8 \textit{attention heads}, dimensi 512), dan \textit{CTC Output Layer} untuk prediksi urutan (95 karakter). Total 27.86M parameter dibekukan selama pelatihan GAN untuk mempertahankan stabilitas dan mencegah catastrophic forgetting. Justifikasi pemilihan arsitektur recognizer dibahas pada Bab II.3.4; recognizer dilatih pada dokumen ANRI era kolonial abad 16--18.

\subsubsection{Diskriminator Dual-Modal Enhanced}
\label{subsubsec:diskriminator-dual-modal}

Arsitektur diskriminator dual-modal merupakan inovasi kunci yang membedakan penelitian ini dari metode restorasi dokumen konvensional. Berbeda dari diskriminator \textit{single-modal} yang hanya mengevaluasi kualitas visual (seperti PatchGAN pada DE-GAN dan ERB-MultiTask), arsitektur dual-modal yang diusulkan mengevaluasi citra dari dua perspektif komplementer secara simultan: (1) perspektif spasial (visual) melalui jalur CNN yang menilai kualitas goresan, tekstur, dan struktur spasial, dan (2) perspektif sekuensial (teks) melalui jalur BiLSTM yang menilai koherensi karakteristik teks dan kelayakan urutan karakter. Integrasi kedua modalitas melalui mekanisme \textit{bilateral cross-modal attention} memungkinkan diskriminator memberikan sinyal adversarial yang lebih informatif: gradien visual mendorong preservasi detail goresan halus, sementara gradien tekstual mendorong konsistensi struktur karakter yang ramah HTR. Arsitektur ini dirancang untuk memberikan umpan balik adversarial yang lebih kaya dengan mempertimbangkan baik kualitas visual maupun koherensi teks, menghasilkan citra restorasi yang optimal untuk dual-objective: realisme visual tinggi dan keterbacaan HTR yang ditingkatkan.

\paragraph{Spesifikasi Arsitektur Diskriminator}

Tabel~\ref{tab:discriminator_arch} merinci arsitektur lengkap Diskriminator. Diskriminator menerima citra ($128 \times 1024 \times 1$) dan sekuens teks (panjang maksimal 128 karakter), kemudian memprosesnya melalui dua jalur paralel seperti ditunjukkan pada Gambar~\ref{fig:discriminator_architecture}.

% Tabel Arsitektur Diskriminator
\begin{table}[H]
\caption[Arsitektur Diskriminator Dual-Modal Enhanced dengan spesifikasi]{Arsitektur Diskriminator Dual-Modal Enhanced. Spesifikasi detail tiga cabang (Citra CNN, Teks BiLSTM, Fusi \textit{Cross-Modal}). Cabang citra memproses fitur visual melalui 4 \textit{DownBlocks} dengan \textit{spatial attention}, cabang teks memproses sekuens karakter melalui embedding + BiLSTM + \textit{self-attention}. \textit{Cross-modal fusion} menghasilkan klasifikasi \textit{real}/\textit{fake}. Total: 17.4M parameter.}
\label{tab:discriminator_arch}
\scriptsize
\centering
\begin{tabular}{|l|l|p{4cm}|r|}
\hline
\textbf{Cabang} & \textbf{Lapisan} & \textbf{Keluaran} & \textbf{Param.} \\
\hline
\multirow{9}{*}{\rotatebox{90}{\tiny CNN}}
& \textit{Input} & $128{\times}1024{\times}1$ & - \\
& Conv2D+BN & $128{\times}1024{\times}64$ & 0.6K \\
& DownBlk-1 & $64{\times}512{\times}64$ & 37K \\
& DownBlk-2 & $32{\times}256{\times}128$ & 148K \\
& DownBlk-3 & $16{\times}128{\times}256$ & 590K \\
& DownBlk-4 & $8{\times}64{\times}512$ & 2.4M \\
& SpatialAttn & $8{\times}64{\times}512$ & 4.7K \\
& ResBlk-512 & $8{\times}64{\times}512$ & 2.4M \\
& AvgPool & $512$ & - \\
\hline
\multirow{5}{*}{\rotatebox{90}{\tiny LSTM}}
& \textit{Input} & $128$ (seq.) & - \\
& Embed & $128{\times}128$ & 12.8K \\
& BiLSTM & $128{\times}512$ & 1.3M \\
& SelfAttn & $128{\times}512$ & 262K \\
& AvgPool & $512$ & - \\
\hline
\multirow{6}{*}{\rotatebox{90}{\tiny Fusi}}
& Proj-CNN & $128$ & 65.7K \\
& Proj-LSTM & $128$ & 65.7K \\
& CrossAttn & $256$ & 33K \\
& Dense-1+Drop & $512$ & 131K \\
& Dense-2+Drop & $256$ & 131K \\
& Sigmoid & $1$ & 257 \\
\hline
\multicolumn{3}{|r|}{\textbf{Total:}} & \textbf{17.4M} \\
\hline
\end{tabular}
\end{table}

% Diagram Arsitektur Diskriminator Dual-Modal
\begin{figure}[H]
\centering
\includegraphics[width=0.5\textwidth,keepaspectratio]{/home/lambda_one/tesis/GAN-HTR-ORI/docRestoration/Paper/drawio/arsitektur-dualModal.png}
\caption[Visualisasi arsitektur Diskriminator Dual-Modal Enhanced]{Arsitektur Diskriminator Dual-Modal Enhanced. Visualisasi dua cabang pemrosesan paralel dengan fusi cross-modal bilateral. Cabang CNN (4 DownBlocks + spatial attention) menghasilkan representasi visual 512-D, cabang BiLSTM (embedding + self-attention) menghasilkan representasi tekstual 512-D. Mekanisme bilateral cross-modal attention (Img$\leftrightarrow$Text) memungkinkan interaksi dua arah untuk validasi holistik koherensi visual-tekstual, menghasilkan klasifikasi real/fake pada lapisan fusi 128-D.}
\label{fig:discriminator_architecture}
\end{figure}

Gambar~\ref{fig:discriminator_architecture} memvisualisasikan arsitektur dual-modal yang membedakan pendekatan ini dari Diskriminator konvensional. Berbeda dari PatchGAN~\cite{isola2017} yang hanya mengevaluasi patch visual, arsitektur ini menggabungkan analisis spasial (CNN) dan sekuensial (BiLSTM) untuk validasi holistik. Karakteristik utama adalah \textit{bilateral cross-modal attention} (Img$\leftrightarrow$Text) yang memungkinkan interaksi dua arah: cabang citra memvalidasi koherensi teks dari HTR recognizer, sementara cabang teks memverifikasi kualitas visual goresan. Mekanisme bilateral ini menghasilkan gradient feedback yang lebih kaya kepada Generator dengan reduksi parameter 52\% (17.4M vs 137M baseline), mendorong perbaikan simultan pada kualitas visual dan konsistensi tekstual.

\paragraph{Cabang Citra (CNN dengan Atensi Spasial)}

Cabang ini mengekstrak fitur visual spasial melalui arsitektur berbasis ResNet. Masukan citra $I \in \mathbb{R}^{128 \times 1024 \times 1}$ diproses melalui lapisan konvolusional awal (64 kanal) dengan \textit{batch normalization} dan aktivasi LeakyReLU ($\alpha=0.2$). Empat blok downsampling residual secara progresif mengurangi resolusi spasial sambil meningkatkan kedalaman kanal:

\begin{equation}
\begin{split}
128 \times 1024 \times 64 &\rightarrow 64 \times 512 \times 64 \rightarrow 32 \times 256 \times 128 \\
&\rightarrow 16 \times 128 \times 256 \rightarrow 8 \times 64 \times 512
\end{split}
\label{eq:IV.4}
\end{equation}

Mekanisme atensi spasial dengan kernel $3 \times 3$ diterapkan pada fitur 512 kanal untuk fokus pada wilayah pembawa teks:

\begin{equation}
\alpha_{\text{spatial}} = \sigma(\text{Conv}_{3 \times 3}(F_{\text{img}}))
\label{eq:IV.5}
\end{equation}

\noindent di mana $\sigma$ adalah fungsi sigmoid dan $F_{\text{img}}$ adalah peta fitur dengan dimensi $8 \times 64 \times 512$. Setelah atensi spasial, blok residual tambahan dengan 512 kanal (ResBlock-512) digunakan untuk pemrosesan fitur lebih lanjut tanpa mengubah dimensi spasial, kemudian global average pooling mengagregasi fitur spasial $8 \times 64 \times 512$ menjadi vektor representasi $F_{\text{img}}^{\text{global}} \in \mathbb{R}^{512}$.

\paragraph{Cabang Teks (BiLSTM dengan Atensi-Mandiri)}

Cabang ini menangkap koherensi sekuensial teks. Sekuens karakter masukan $S \in \mathbb{N}^{L}$ (panjang $L \leq 128$) pertama kali diproyeksikan ke ruang embedding $E \in \mathbb{R}^{L \times 128}$. BiLSTM memproses sekuens dalam kedua arah untuk menangkap dependensi kontekstual:

\begin{equation}
h_t = [\overrightarrow{h_t}; \overleftarrow{h_t}] = [\text{LSTM}_{\text{fwd}}(E_t); \text{LSTM}_{\text{bwd}}(E_t)]
\label{eq:IV.6}
\end{equation}

\noindent menghasilkan \textit{hidden states} $H \in \mathbb{R}^{L \times 512}$. Mekanisme atensi-mandiri \textit{multi-head scaled dot-product} diterapkan:

\begin{equation}
\text{Attention}(Q, K, V) = \text{softmax}\left(\frac{QK^T}{\sqrt{d_k}}\right)V
\label{eq:IV.7}
\end{equation}

\noindent di mana $Q, K, V$ adalah proyeksi linear dari $H$. \textit{Global average pooling} menghasilkan vektor teks $F_{\text{text}}^{\text{global}} \in \mathbb{R}^{512}$.

\paragraph{Fusi Cross-Modal}

Kedua vektor fitur $F_{\text{img}}^{\text{global}} \in \mathbb{R}^{512}$ dan $F_{\text{text}}^{\text{global}} \in \mathbb{R}^{512}$ diproses melalui mekanisme atensi cross-modal bilateral untuk memungkinkan interaksi dua arah antar modalitas. Pendekatan bilateral diperlukan karena: (1) citra perlu memvalidasi koherensi teks yang diprediksi pengenal HTR, dan (2) teks perlu memverifikasi kualitas visual goresan yang menghasilkan karakter tersebut. Dengan demikian, kedua modalitas saling memberikan konteks untuk deteksi anomali yang lebih robust. Setiap modalitas memproyeksikan fiturnya ke ruang bersama berdimensi rendah $d=128$ melalui enam matriks proyeksi terpisah:

Atensi Citra$\rightarrow$Teks: Citra sebagai \textit{query}, teks sebagai \textit{key}/\textit{value}:
\begin{align}
Q_{\text{img}} &= W_{\text{img}}^Q F_{\text{img}}^{\text{global}} \nonumber\\
K_{\text{text}} &= W_{\text{text}}^K F_{\text{text}}^{\text{global}} \nonumber\\
V_{\text{text}} &= W_{\text{text}}^V F_{\text{text}}^{\text{global}} \label{eq:IV.8}\\
F_{\text{img}}^{\text{att}} &= \text{softmax}\left(\frac{Q_{\text{img}} K_{\text{text}}^T}{\sqrt{d}}\right) V_{\text{text}} \label{eq:IV.9}
\end{align}

\noindent di mana $W_{\text{img}}^Q, W_{\text{text}}^K, W_{\text{text}}^V \in \mathbb{R}^{512 \times 128}$ adalah matriks proyeksi yang dipelajari.

Atensi Teks$\rightarrow$Citra: Teks sebagai \textit{query}, citra sebagai \textit{key}/\textit{value}:
\begin{align}
Q_{\text{text}} &= W_{\text{text}}^Q F_{\text{text}}^{\text{global}} \nonumber\\
K_{\text{img}} &= W_{\text{img}}^K F_{\text{img}}^{\text{global}} \nonumber\\
V_{\text{img}} &= W_{\text{img}}^V F_{\text{img}}^{\text{global}} \label{eq:IV.10}\\
F_{\text{text}}^{\text{att}} &= \text{softmax}\left(\frac{Q_{\text{text}} K_{\text{img}}^T}{\sqrt{d}}\right) V_{\text{img}} \label{eq:IV.11}
\end{align}

\noindent di mana $W_{\text{text}}^Q, W_{\text{img}}^K, W_{\text{img}}^V \in \mathbb{R}^{512 \times 128}$ adalah matriks proyeksi simetris.

Kedua fitur yang telah mengalami atensi bilateral ($F_{\text{img}}^{\text{att}}, F_{\text{text}}^{\text{att}} \in \mathbb{R}^{128}$) digabungkan menjadi vektor fusi: $F_{\text{fused}} = [F_{\text{img}}^{\text{att}}; F_{\text{text}}^{\text{att}}] \in \mathbb{R}^{256}$. Kepala klasifikasi terdiri dari dua lapisan padat (512 $\rightarrow$ 256 neuron) dengan dropout dan aktivasi LeakyReLU, diikuti lapisan keluaran dengan aktivasi sigmoid yang menghasilkan skor validitas $D(I, S) \in [0, 1]$ untuk membedakan pasangan citra-teks asli dari pasangan palsu yang dihasilkan generator.

\paragraph{Konfigurasi Hyperparameter}

Kernel atensi spasial menggunakan ukuran $3 \times 3$; dimensi proyeksi cross-modal ditetapkan 128; unit LSTM per arah sebanyak 256 menghasilkan total keluaran BiLSTM 512 dimensi. \textit{Batch normalization momentum}~$=0.9$; tingkat dropout~$=0.1$ pada kepala klasifikasi dan dropout~$=0.2$ pada embedding teks; fungsi aktivasi LeakyReLU ($\alpha=0.2$) untuk CNN dan sigmoid untuk aktivasi rekuren LSTM. Pemilihan kernel $3\times3$ mengurangi artefak pada goresan tipis dibanding kernel lebih besar, sementara dropout rendah mempertahankan informasi detail yang krusial untuk diskriminasi dokumen terdegradasi.

\paragraph{Motivasi dan Rasional Desain}

Pendekatan dual-modal ini mempertimbangkan bahwa dokumen teks memiliki sifat ganda: (1) penampilan visual (kualitas goresan, tingkat derau), (2) struktur sekuensial (aliran karakter, jarak antar kata). Jalur CNN menangkap kualitas spasial, sementara BiLSTM menangkap koherensi sekuensial. Fusi cross-modal bilateral memungkinkan diskriminator menilai pasangan citra-teks secara holistik dengan interaksi dua arah, memberikan sinyal adversarial yang lebih informatif kepada generator untuk perbaikan simultan pada kualitas visual dan konsistensi tekstual.

Diskriminator yang diusulkan mencapai efisiensi parameter signifikan (17.4M, reduksi 52\% dari baseline 137M) melalui arsitektur yang disederhanakan tanpa mengorbankan kemampuan diskriminatif: (1) proyeksi dimensi rendah (128) untuk fusi cross-modal mengurangi kompleksitas komputasi sambil mempertahankan kapasitas representasi; (2) kernel atensi spasial $3\times3$ fokus pada detail lokal; (3) blok residual yang dioptimalkan menyeimbangkan kedalaman arsitektur dengan stabilitas pelatihan. Kombinasi \textit{batch normalization momentum}~$=0.9$ dan dropout~$=0.1$ menghasilkan statistik stabil dengan pelestarian informasi goresan tipis.

Dengan kombinasi analisis spasial CNN dan pemrosesan sekuens BiLSTM, diskriminator dapat mendeteksi:
\begin{itemize}
    \item Goresan yang rusak atau terputus yang mengganggu alur teks (melalui deteksi anomali sekuens BiLSTM)
    \item Pola derau, artefak, dan tembus-balik (melalui fitur spasial CNN)
    \item Lebar goresan yang tidak konsisten (melalui mekanisme atensi spasial)
    \item Celah atau koneksi yang tidak wajar antar karakter (melalui analisis fusi cross-modal)
\end{itemize}

Dengan menggabungkan penilaian spasial dan sekuensial, diskriminator memberikan umpan balik yang lebih kaya kepada generator, mendorong produksi citra yang bersih secara visual dan koheren secara tekstual. Efisiensi parameter yang dicapai (17.4M, reduksi 52\% dari baseline) membuktikan bahwa arsitektur dual-modal yang diusulkan tidak memerlukan kompleksitas berlebihan untuk mencapai kemampuan diskriminatif yang efektif---fusi \textit{bilateral cross-attention} dengan dimensi proyeksi rendah (128-D) menghasilkan interaksi modalitas yang kaya tanpa overhead komputasional signifikan.

Mode input teks diskriminator:
Diskriminator dual-modal dapat beroperasi dalam dua mode input teks: (1) Ground Truth Mode menggunakan label teks yang sudah dianotasi, memberikan sinyal pembelajaran yang bersih namun kurang realistis untuk deployment, dan (2) Predicted Mode menggunakan prediksi dari frozen recognizer, lebih realistis namun rentan terhadap error propagation. Perbandingan empiris kedua mode dan justifikasi pemilihan mode untuk training diuraikan pada Bab V (Evaluasi dan Analisis).

\vspace{0.5em}
\noindent\textit{Rancangan arsitektur yang telah diuraikan pada Subsection~\ref{subsec:rancangan-arsitektur} (IV.1.1--IV.1.6) memberikan fondasi teoretis untuk framework GAN-HTR yang diusulkan, mencakup spesifikasi generator U-Net Enhanced (21.8M parameter), frozen recognizer hibrida CNN-Transformer (27.86M parameter), diskriminator dual-modal dengan fusi cross-modal bilateral (17.4M parameter), dan fungsi loss multi-komponen yang dioptimalkan. Implementasi praktis rancangan ini memerlukan infrastruktur komputasi yang memadai dan konfigurasi eksperimental yang teliti untuk memastikan stabilitas pelatihan dan reproduktibilitas hasil. Subsection berikut menguraikan lingkungan eksperimen, pengaturan dataset, protokol training, dan metrik evaluasi yang digunakan untuk validasi empiris kerangka kerja yang diusulkan.}
\vspace{0.5em}

\subsubsection{Fungsi Loss Multikomponen}
\label{subsubsec:fungsi-loss-multikomponen}

Generator dilatih dengan kombinasi lima komponen loss yang dikalibrasi melalui penyetelan empiris berdasarkan analisis inverse scaling terhadap distribusi magnitudo loss dan divalidasi secara post-hoc menggunakan pencarian kisi sistematis (27 konfigurasi dengan metode orthogonal array sampling) yang mengonfirmasi optimalitas konfigurasi yang dipilih. Fungsi loss total menggabungkan komponen-komponen tersebut secara terbobot untuk mengoptimalkan kualitas visual dan keterbacaan teks secara simultan:

\begin{equation}
\begin{split}
\mathcal{L}_{\text{total}} = &\lambda_{\text{adv}}\mathcal{L}_{\text{adv}} + \lambda_{\text{pixel}}\mathcal{L}_{\text{pixel}} + \lambda_{\text{perc}}\mathcal{L}_{\text{perc}} \\
&+ \lambda_{\text{ctc}}\mathcal{L}_{\text{ctc}} + \lambda_{\text{rec-feat}}\mathcal{L}_{\text{rec-feat}}
\end{split}
\label{eq:IV.12}
\end{equation}

Konfigurasi bobot loss untuk setiap komponen ditunjukkan dalam Tabel~\ref{table:optimal_weights}. Studi ablasi sistematis (Bab~V) mengidentifikasi konfigurasi optimal 4-komponen (tanpa $\mathcal{L}_{\text{rec-feat}}$) sebagai keseimbangan \textit{Pareto} antara kualitas visual dan keterbacaan HTR.

\paragraph{Adversarial Loss}

\textit{Adversarial loss} standar GAN mendorong pembuatan citra yang realistis dengan melatih generator untuk menipu diskriminator:

\begin{equation}
\mathcal{L}_{\text{adv}} = -\log D(I_{\text{deg}}, I_{\text{gen}})
\label{eq:IV.13}
\end{equation}

di mana $D(I_{\text{deg}}, I_{\text{gen}})$ adalah probabilitas diskriminator menganggap pasangan citra terdegradasi masukan $I_{\text{deg}}$ dan citra hasil restorasi $I_{\text{gen}}$ sebagai pasangan autentik. Diskriminator dilatih untuk memaksimalkan kemampuan membedakan:

\begin{equation}
\begin{split}
\mathcal{L}_D = -[&\log D(I_{\text{deg}}, I_{\text{gt}}) \\
&+ \log(1 - D(I_{\text{deg}}, I_{\text{gen}}))]
\end{split}
\label{eq:IV.14}
\end{equation}

di mana pasangan $(I_{\text{deg}}, I_{\text{gt}})$ merupakan pasangan nyata (citra terdegradasi dengan ground truth bersih), sedangkan $(I_{\text{deg}}, I_{\text{gen}})$ merupakan pasangan palsu hasil generator. Mekanisme adversarial ini mendorong generator menghasilkan citra dengan distribusi tekstur dan karakteristik visual yang tidak dapat dibedakan dari dokumen bersih autentik.

\paragraph{Pixel Reconstruction Loss}

\textit{Pixel reconstruction loss} menggunakan norma L1 untuk memastikan kesamaan piksel-demi-piksel dengan \textit{ground truth}:

\begin{equation}
\mathcal{L}_{\text{pixel}} = \|I_{\text{gen}} - I_{\text{gt}}\|_1
\label{eq:IV.15}
\end{equation}

Norma L1 dipilih karena menghasilkan hasil yang lebih tajam dibandingkan norma L2 dan lebih robust terhadap pencilan. Berbeda dengan L2 yang memberikan penalti kuadratik (cenderung menghasilkan citra kabur), L1 memberikan penalti linear yang mempertahankan ketajaman tepi dan detail goresan halus yang penting untuk keterbacaan HTR.

\paragraph{Perceptual Loss}

\textit{Perceptual loss} \cite{johnson2016} membandingkan fitur tingkat tinggi dari jaringan VGG-19 yang telah dilatih pada \textit{ImageNet}:

\begin{equation}
\mathcal{L}_{\text{perc}} = \sum_{l \in L} \|\phi_l(I_{\text{gen}}) - \phi_l(I_{\text{gt}})\|_2
\label{eq:IV.16}
\end{equation}

di mana $\phi_l$ menunjukkan fitur dari lapisan $l$ dari jaringan VGG-19 praterlatih, dan $L$ adalah himpunan lapisan konvolusional yang dipilih untuk ekstraksi fitur. Fitur diekstrak dari beberapa lapisan konvolusional untuk perbandingan perseptual multiskala. Pendekatan ini memastikan pelestarian karakteristik visual dari detail tingkat rendah (tepi, tekstur) hingga fitur semantik tingkat tinggi (pola goresan, bentuk karakter), menghasilkan citra yang lebih koheren secara perseptual meskipun terdapat perbedaan piksel minor.

\paragraph{CTC Loss}

CTC loss memberikan gradien sadar-teks langsung dari frozen recognizer untuk mengoptimalkan keterbacaan HTR secara eksplisit. Connectionist Temporal Classification (CTC) dipilih karena kemampuannya menangani sekuens dengan panjang variabel tanpa memerlukan alignment eksplisit antara input citra dan output teks, ideal untuk dokumen tulisan tangan dengan variasi spasial karakter yang tinggi. Formulasi matematis CTC dan implementasi teknis telah dijelaskan secara rinci pada Subbab~\ref{subsubsec:ekstraksi-fitur-recognizer}. Komponen ini merupakan inovasi kunci penelitian yang membedakan kerangka kerja yang diusulkan dari metode restorasi dokumen konvensional yang hanya mengoptimalkan kualitas visual tanpa mempertimbangkan keterbacaan mesin. Analisis kontribusi menunjukkan CTC loss berkontribusi 67.5\% terhadap total generator loss (meskipun bobot $\lambda_{\text{ctc}}=0.15$ relatif kecil) karena magnitudo \textit{natural} CTC ($\sim$100--300) jauh lebih besar dibanding komponen visual ($\sim$0.1--5.0), memvalidasi peran dominan komponen HTR dalam optimasi. Strategi \textit{clipping} dengan batas maksimal 400 diterapkan untuk mencegah gradien ekstrem pada kasus degradasi berat yang dapat mengganggu stabilitas pelatihan adversarial.

\paragraph{Recognition Feature Loss}

Komponen \textit{recognition feature loss} mendorong karakteristik yang ramah HTR dengan meminimalkan jarak fitur dari \textit{encoder recognizer} beku. Formulasi matematis dan justifikasi telah dijelaskan pada Subbab~\ref{subsubsec:ekstraksi-fitur-recognizer}. Komponen ini divalidasi dalam studi ablasi (Bab~V) dan \textit{grid search post-hoc} yang menguji RecFeat$\in\{0, 8, 16\}$, namun hasil menunjukkan kontribusi marginal (0.1\% dari total loss) dan korelasi tinggi dengan CTC loss (Pearson $r=0.89$). Berdasarkan analisis multi-objektif, komponen ini dinonaktifkan pada konfigurasi produksi final ($\lambda_{\text{rec-feat}}=0$), memvalidasi bahwa \textit{komponen HTR tunggal (CTC loss) mencukupi} untuk optimasi berorientasi HTR asalkan diberi bobot yang tepat berdasarkan \textit{magnitude natural}. Keputusan ini merupakan \textit{hasil proses optimasi}, bukan keterbatasan metode, dan meningkatkan efisiensi komputasi tanpa mengorbankan performa HTR.

\begin{table}[H]
\renewcommand{\arraystretch}{1.3}
\caption[Konfigurasi bobot loss optimal dengan prinsip inverse scaling]{Konfigurasi bobot loss optimal. Bobot dipilih berdasarkan prinsip inverse scaling empiris\\
(w$_i$ $\propto$ 1/magnitude$_i$) dari analisis distribusi loss, kemudian divalidasi melalui pencarian kisi post-hoc\\
(27 konfigurasi) yang mengonfirmasi konfigurasi ini optimal pada validation set. Bobot pixel L1 yang tinggi ($\lambda=50$)\\
memastikan kesetiaan rekonstruksi dasar, bobot CTC yang relatif rendah ($\lambda=0.15$) dikompensasi oleh magnitudo natural loss CTC\\
($\sim$100--300) dibandingkan loss L1 ($\sim$0.1--0.5). Validasi grid search menunjukkan tidak ada konfigurasi alternatif\\
yang memberikan peningkatan signifikan pada fase awal training (3 epoch).}
\label{table:optimal_weights}
\small
\centering
\begin{tabular}{|l|p{3cm}|p{6cm}|}
\hline
\textbf{Komponen Loss} & \textbf{Bobot ($\lambda$)} & \textbf{Justifikasi} \\
\hline
\textit{adversarial} ($\mathcal{L}_{\text{adv}}$) & 3.0 & Realisme tekstur dan distribusi visual yang ditingkatkan tanpa mengorbankan stabilitas pelatihan \\
\hline
Piksel L1 ($\mathcal{L}_{\text{pixel}}$) & 50.0 & Pelestarian struktur piksel yang seimbang; bobot tinggi memastikan kesetiaan rekonstruksi dasar \\
\hline
\textit{Perceptual} ($\mathcal{L}_{\text{perc}}$) & 1.0 & Pelestarian topologi dan koherensi fitur semantik tingkat tinggi \\
\hline
Recognition Feature loss ($\mathcal{L}_{\text{rec-feat}}$) & 8.0 & Panduan karakteristik ramah HTR yang kuat (dinonaktifkan produksi: $\lambda=0$) \\
\hline
CTC loss ($\mathcal{L}_{\text{ctc}}$) & 0.15 & Pemantauan keterbacaan HTR dengan \textit{clipping} numerik untuk stabilitas gradien \\
\hline
\end{tabular}
\end{table}

Pemilihan bobot loss (Tabel~\ref{table:optimal_weights}) dilakukan berdasarkan prinsip inverse scaling empiris dari analisis distribusi magnitudo loss, dengan pertimbangan bahwa komponen dengan magnitudo natural lebih besar memerlukan bobot lebih kecil untuk mencegah dominasi berlebihan dalam total generator loss. Konfigurasi yang dipilih (pixel=50.0, adversarial=3.0, perceptual=1.0, CTC=0.15, recognition feature=8.0) kemudian divalidasi melalui dua mekanisme: (1) validasi post-hoc menggunakan pencarian kisi sistematis dengan 27 konfigurasi (pixel$\in\{25,50,100\}$, adversarial$\in\{1.5,3.0,6.0\}$, perceptual$\in\{0.5,1.0,2.0\}$, CTC$\in\{0.075,0.15,0.30\}$, recognition feature$\in\{0,8,16\}$) yang menggunakan metode orthogonal array sampling, menunjukkan bahwa konfigurasi production berada dalam rentang optimal dan tidak ada konfigurasi alternatif yang memberikan peningkatan signifikan pada fase awal konvergensi (3 epoch); (2) validasi GradNorm selama pelatihan production (50 epoch) menunjukkan variasi bobot 0\%, mengonfirmasi konfigurasi sudah berada pada titik keseimbangan optimal untuk konvergensi jangka panjang.

Optimasi Berorientasi HTR: Istilah ``optimasi fungsi loss berorientasi HTR'' pada penelitian ini mengacu pada strategi kalibrasi yang \textit{menempatkan komponen HTR (CTC loss) sebagai faktor dominan} dalam pemilihan bobot, bukan sekadar komponen tambahan pada fungsi loss visual. Meskipun bobot CTC tampak rendah secara nominal ($\lambda_{\text{ctc}}=0.15$), magnitudo natural loss CTC yang besar (rerata 392.01, di-\textit{clip} ke maksimum 400.0) menghasilkan kontribusi efektif 67.5\% terhadap total \textit{generator loss} (analisis detail pada Bab~V). Sebagai perbandingan, komponen visual (pixel + adversarial + perceptual) hanya berkontribusi 32.4\% secara kumulatif. Dominasi CTC loss ini memastikan generator \textit{secara eksplisit dioptimalkan untuk keterbacaan HTR}, bukan hanya kualitas visual, membedakan pendekatan ini dari metode restorasi dokumen konvensional yang tidak mempertimbangkan metrik keterbacaan teks mesin dalam fungsi objektif.

Karakteristik dataset menunjukkan dominasi goresan tipis yang memerlukan keseimbangan khusus antara rekonstruksi piksel dan panduan berbasis teks. Bobot rekonstruksi piksel yang tinggi ($\lambda_{\text{pixel}}=50.0$) memberikan sinyal rekonstruksi kuat dengan tetap memberikan fleksibilitas untuk peningkatan adversarial, sementara bobot CTC yang relatif rendah ($\lambda_{\text{ctc}}=0.15$) dikompensasi oleh magnitudo natural loss CTC yang besar ($\sim$100--300) dibandingkan loss L1 ($\sim$0.1--0.5).

Analisis \textit{Trade-off} Multi-Objektif: Orientasi HTR pada optimasi loss dibuktikan melalui analisis \textit{trade-off} antara metrik kualitas visual (PSNR, SSIM) dan metrik keterbacaan teks (CER, WER) pada validation set. Validasi \textit{grid search post-hoc} (detail pada Bab~V) menunjukkan bahwa peningkatan bobot CTC ($0.15 \rightarrow 0.30$) meningkatkan kontribusi HTR tetapi menurunkan PSNR pada fase awal pelatihan, sementara penurunan bobot CTC ($0.15 \rightarrow 0.075$) meningkatkan PSNR tetapi mengorbankan keterbacaan teks. Konfigurasi $\lambda_{\text{ctc}}=0.15$ dipilih sebagai titik keseimbangan \textit{Pareto} yang mempertahankan kualitas visual tinggi (PSNR 30.91 dB pada best checkpoint epoch 44) sambil mencapai peningkatan keterbacaan HTR signifikan dibandingkan metode \textit{baseline} (CER 34.9\% vs 64.2\% pada citra terdegradasi, reduksi relatif 45.7\%). Analisis lengkap \textit{Pareto front} dan sensitivitas komponen tersedia pada Bab~V, dengan validasi empiris bahwa konfigurasi yang dipilih berada pada rentang optimal untuk objektif ganda (visual + HTR).

\vspace{0.5em}
\noindent\textit{Rancangan arsitektur yang telah diuraikan pada Subsection~\ref{subsec:rancangan-arsitektur} (IV.1.1--IV.1.6) memberikan fondasi teoretis untuk framework GAN-HTR yang diusulkan, mencakup spesifikasi generator U-Net Enhanced (21.8M parameter), frozen recognizer hibrida CNN-Transformer (27.86M parameter), diskriminator dual-modal dengan fusi cross-modal bilateral (17.4M parameter), dan fungsi loss multi-komponen yang dioptimalkan. Implementasi praktis rancangan ini memerlukan infrastruktur komputasi yang memadai dan konfigurasi eksperimental yang teliti untuk memastikan stabilitas pelatihan dan reproduktibilitas hasil. Subsection berikut menguraikan lingkungan eksperimen, pengaturan dataset, protokol training, dan metrik evaluasi yang digunakan untuk validasi empiris kerangka kerja yang diusulkan.}
\vspace{0.5em}

\subsection{Lingkungan Eksperimen dan Spesifikasi Komputasi}
\label{subsec:lingkungan-eksperimen}

Implementasi menggunakan \textit{workstation high-performance} dengan spesifikasi yang dirinci pada Tabel~\ref{tab:hardware-specs-actual}. Sistem ini menyediakan komputasi paralel GPU yang diperlukan untuk pelatihan efisien arsitektur GAN kompleks dengan komponen dual-modal dan frozen recognizer.

\begin{table}[H]
\caption[Spesifikasi lingkungan eksperimen dengan detail hardware]{Spesifikasi lingkungan eksperimen. Detail hardware (GPU RTX A4000 16GB, CPU AMD Ryzen Threadripper PRO 3955WX, 128GB RAM), sistem operasi (Ubuntu 22.04.5 LTS), dan software framework yang digunakan untuk semua eksperimen pelatihan dan evaluasi dalam penelitian ini.}
\label{tab:hardware-specs-actual}
\small
\centering
\begin{tabular}{|l|l|}
\hline
\textbf{Komponen} & \textbf{Spesifikasi} \\
\hline
GPU & 2$\times$ NVIDIA RTX A4000 (16 GB GDDR6 per GPU) \\
\hline
CPU & AMD Ryzen Threadripper PRO 3955WX (16-core/32-thread) \\
\hline
RAM & 128 GB DDR4-3200 \\
\hline
Storage & 1 TB NVMe SSD (Samsung 980 PRO) \\
\hline
OS & Ubuntu 22.04.5 LTS \\
\hline
Framework & TensorFlow 2.16.1 + CUDA 12.8 + cuDNN 8.9.7 \\
\hline
\end{tabular}
\end{table}

\paragraph{Infrastruktur Software dan Protokol Reproduktibilitas}

Implementasi sistem menggunakan \textit{software stack} terintegrasi yang dipilih untuk memastikan stabilitas numerik, efisiensi komputasi, dan reproduktibilitas penuh. Konfigurasi utama mencakup Python 3.10.12 sebagai runtime environment, TensorFlow 2.16.1 dengan dukungan GPU melalui CUDA 12.8 dan cuDNN 8.9, serta Poetry 1.7.1 untuk manajemen dependensi deterministik dengan version locking eksplisit semua library (file \texttt{poetry.lock} memastikan instalasi identik lintas environment). TensorFlow Text 2.16 digunakan untuk operasi CTC loss, sementara MLflow 2.14.1 mengelola pelacakan eksperimen dengan auto-logging semua hyperparameter, metrik training, dan artefak model.

Implementasi menggunakan presisi numerik FP32 (float32) untuk semua operasi training, dipilih karena stabilitas superior CTC loss pada presisi penuh dibandingkan FP16 mixed precision yang rentan terhadap \textit{numerical underflow} pada probabilitas karakter kecil. Penjadwalan laju pembelajaran adaptif diimplementasikan dengan strategi tiga tahap: (1) \textit{linear warmup} selama 10 epoch pertama untuk stabilitas awal, (2) \textit{exponential decay} dengan faktor 0.95 setiap 5 epoch untuk konvergensi halus, dan (3) callback \texttt{ReduceLROnPlateau} yang memantau validation loss dan mengurangi learning rate 50\% jika tidak ada perbaikan selama 3 epoch berturut-turut. Mekanisme \textit{adaptive loss balancing} memantau rasio kontribusi setiap komponen loss dan melakukan penyesuaian minor ($\pm$5\%) jika terjadi dominasi berlebihan, memastikan keseimbangan dinamis antara objektif visual dan tekstual (detail implementasi pada Lampiran B.2).

Protokol reproduktibilitas dikonfigurasi dengan empat mekanisme deterministik: (1) seed acak tetap (\texttt{seed=42}) untuk NumPy, TensorFlow, dan Python \texttt{random}, (2) mode deterministik TensorFlow (\texttt{tf.config.experimental.enable\_op\_determinism()}) dan cuDNN (\texttt{TF\_DETERMINISTIC\_OPS=1}) untuk memastikan operasi GPU identik antar run, (3) konfigurasi GPU single-device (GPU:0) untuk menghindari variabilitas paralelisme multi-GPU, dan (4) MLflow experiment tracking yang merekam semua konfigurasi, checkpoint model, dan log training untuk audit penuh. Strategi \textit{checkpointing} menyimpan model setiap 5 epoch plus \textit{best model} berdasarkan metrik gabungan PSNR-CER pada validation set. Konfigurasi lengkap environment dan dependency tersedia pada Lampiran A.4 untuk replikasi eksak.

Implikasi konfigurasi:

Infrastruktur komputasi yang dikonfigurasi (2$\times$ RTX A4000, 128 GB RAM) memungkinkan pelatihan framework GAN-HTR dengan ukuran batch yang memadai (2 per GPU) untuk citra beresolusi tinggi ($1024 \times 128$ piksel) tanpa mengorbankan detail goresan halus paleografi. Presisi FP32 yang dipilih memastikan stabilitas numerik CTC loss dengan trade-off waktu pelatihan yang dapat diterima (48 GPU-hours untuk 50 epoch). Konfigurasi reproduktibilitas deterministik memungkinkan replikasi eksperimen dengan hasil identik, krusial untuk validasi studi ablasi sistematis pada Bab V.

\subsection{Pengaturan Eksperimental}
\label{subsec:pengaturan-eksperimental}

Evaluasi komprehensif framework GAN-HTR yang diusulkan memerlukan strategi dataset ganda yang melengkapi keterbatasan inheren evaluasi restorasi dokumen historis. Penelitian ini mengadopsi pendekatan evaluasi dual: (1) dataset semi-sintetis realistis untuk evaluasi kuantitatif presisi dengan metrik objektif (PSNR, SSIM, CER) yang memerlukan pasangan citra terdegradasi-bersih dengan keselarasan piksel sempurna---tidak tersedia pada dokumen historis nyata karena degradasi alami yang ireversibel, dan (2) dataset dokumen historis nyata ANRI untuk validasi generalisasi kualitatif pada degradasi autentik yang kompleks dan tidak dapat direplikasi sepenuhnya melalui augmentasi komputasional. Strategi dual ini mengikuti praktik standar riset restorasi dokumen historis, menyeimbangkan kebutuhan pengukuran kuantitatif akurat dengan validasi performa pada kondisi dunia nyata.

Bagian ini menjelaskan konfigurasi eksperimental komprehensif yang digunakan untuk evaluasi kerangka kerja yang diusulkan, mencakup konstruksi kedua dataset, metode baseline klasik untuk perbandingan kuantitatif, dan protokol pengujian untuk memastikan rigor ilmiah dan reproduktibilitas penuh.

\subsubsection{Dataset Semi-Sintetis Realistis untuk Evaluasi Kuantitatif}

Alasan memilih metode:

Evaluasi kuantitatif restorasi dokumen memerlukan pasangan citra terdegradasi-bersih dengan keselarasan tingkat piksel yang sempurna untuk mengukur metrik PSNR, SSIM, dan CER secara akurat. Dokumen historis nyata yang mengalami degradasi alami tidak memiliki ground truth bersih karena tidak dapat dipulihkan ke kondisi asli---kerusakan fisik dan kimia selama berabad-abad bersifat ireversibel. Oleh karena itu, dataset semi-sintetis realistis dikonstruksi dengan menggabungkan tiga komponen autentik untuk mempertahankan realisme degradasi historis sambil menyediakan ground truth berpasangan yang diperlukan untuk pengukuran presisi.

Komponen 1: Tulisan Tangan Paleografi Autentik (REAL)

Dokumen arsip ANRI autentik abad ke-16 hingga ke-18 yang relatif bersih (minimal degradasi visual) digunakan sebagai sumber tulisan tangan. Dokumen dianotasi dalam format Transkribus PAGE XML, dibinarisasi menggunakan model DE-GAN praterlatih, dan baris teks individual diekstrak menggunakan koordinat baseline dengan padding vertikal, dinormalisasi ke dimensi $128 \times 1024$ piksel. Hasil ekstraksi adalah citra baris teks dengan karakteristik paleografi autentik---variasi ketebalan goresan, ligatur kompleks, gaya penulisan individu penulis historis---yang tidak dapat direplikasi oleh rendering fon sintetis atau model generatif.

Komponen 2: Pola Degradasi Historis Autentik (REAL)

Dokumen ANRI terpisah yang mengalami degradasi alami berat (penuaan kertas, noda organik, foxing, tinta luntur, korosi tinta \textit{iron gall}) digunakan sebagai sumber tekstur latar belakang degradasi. \textit{Patch} latar belakang berukuran $128 \times 1024$ piksel diekstrak dari wilayah non-teks dokumen rusak ini, mempertahankan pola degradasi historis autentik yang terbentuk selama berabad-abad penyimpanan arsip.

Komponen 3: Augmentasi Degradasi Terkontrol (SEMI-SYNTHETIC)

% \noindent\textit{Catatan: Jalur degradasi komputasional enam tahap yang diuraikan berikut ini diimplementasikan dalam pipeline persiapan dataset (Subsection~\ref{subsubsec:synthetic-degradation-pipeline}, IV.5.2) untuk menghasilkan pasangan citra terdegradasi-bersih dengan keselarasan piksel sempurna.}

Untuk meningkatkan variasi dan menciptakan pasangan terdegradasi-bersih dengan keselarasan piksel sempurna, jalur degradasi berlapis dengan enam tahap komputasional diterapkan secara stokastik pada kombinasi teks autentik dan latar belakang degradasi nyata:

Tahap 1 --- \textit{Overlay} Latar Belakang Tekstur:

Citra teks bersih dikomposisikan pada patch tekstur dokumen nyata yang diekstrak dari koleksi dokumen rusak ANRI. Wilayah teks diredup untuk mensimulasikan penyerapan tinta ke dalam kertas yang sudah tua:
\begin{equation}
I_{\text{degradasi}} = I_{\text{bg}} \cdot (1 - M_{\text{teks}}) + (\alpha \cdot I_{\text{bg}}) \cdot M_{\text{teks}}
\label{eq:overlay-texture}
\end{equation}
dengan $M_{\text{teks}}$ adalah masker teks yang dinormalisasi dan $\alpha$ adalah faktor peredupan.

Tahap 2 --- Efek \textit{Bleed-Through}:

Mensimulasikan tinta yang menembus dari sisi sebalik dokumen dengan membalik citra teks acak secara horizontal, menerapkan \textit{Gaussian blur}, dan mengurangi dari latar belakang.

Tahap 3 --- Noda Organik Berbasis \textit{Perlin Noise}:

Menghasilkan pola noda yang tampak alami menggunakan \textit{Perlin noise} 3D dengan parameter yang divariasikan untuk membuat pola noda organik yang halus dan realistis.

Tahap 4 --- \textit{Foxing Spots} (Bintik Penuaan):

Mensimulasikan perubahan warna terkait usia dengan menambahkan bintik kecil berwarna kecoklatan dengan intensitas acak, mereplikasi pola degradasi biologis pada dokumen historis.

Tahap 5 --- Kerusakan Fisik:

Mensimulasikan robekan dan lipatan dokumen melalui garis acak dan patch yang digelapkan, menciptakan artefak struktural yang konsisten dengan dokumen berusia ratusan tahun.

Tahap 6 --- \textit{Blur} Acak:

Menerapkan \textit{Gaussian blur} (degradasi halus), \textit{median blur} (penghilangan derau), atau \textit{motion blur} (artefak pemindai) dengan probabilitas yang sama untuk mensimulasikan berbagai kondisi akuisisi dan penyimpanan.

Efek-efek ini diterapkan secara berurutan dengan kombinasi yang diacak, menghasilkan dokumen terdegradasi yang beragam dan realistis sambil mempertahankan keselarasan \textit{ground truth} tingkat piksel yang sempurna, menghilangkan kesalahan anotasi manual yang melekat pada dataset dokumen nyata.

Statistik dataset semi-sintetis realistis:

Dataset terdiri dari 4.739 pasangan citra terdegradasi-bersih ($128 \times 1024$ piksel, grayscale, dinormalisasi [0,1]) yang dikonstruksi dari dokumen ANRI autentik. Setiap pasangan menggabungkan: (1) tulisan tangan paleografi historis nyata, (2) pola degradasi dari arsip autentik yang mengalami kerusakan alami, dan (3) augmentasi terkontrol untuk variasi. Dataset ini mempertahankan karakteristik autentik dokumen historis sambil menyediakan ground truth berpasangan yang diperlukan untuk evaluasi kuantitatif presisi tinggi.

\begin{figure}[H]
\centering

\includegraphics[width=0.95\textwidth]{/home/lambda_one/tesis/GAN-HTR-ORI/docRestoration/Paper/data_dukung/gambar_degGT/gt1.png}
\vspace{0.2cm}

\includegraphics[width=0.95\textwidth]{/home/lambda_one/tesis/GAN-HTR-ORI/docRestoration/Paper/data_dukung/gambar_degGT/deg1.png}
\caption[Pasangan dataset semi-sintetis realistis]{Pasangan dataset semi-sintetis realistis. (Atas) Citra ground truth bersih dari dokumen ANRI abad ke-16--18 menunjukkan tulisan tangan paleografi autentik dengan variasi goresan, ligatur kompleks, dan gaya penulisan historis yang tidak dapat direplikasi oleh fon sintetis. (Bawah) Citra terdegradasi menggabungkan pola degradasi historis nyata (penuaan kertas, noda organik dari arsip ANRI autentik) dengan augmentasi terkontrol enam tahap (overlay tekstur, bleed-through, noda \textit{Perlin noise}, \textit{foxing spots}, kerusakan fisik, \textit{blur} acak), mempertahankan realisme degradasi dokumen berusia berabad-abad. Keselarasan piksel sempurna antara pasangan memungkinkan evaluasi kuantitatif akurat dengan metrik PSNR, SSIM, dan CER.}
\label{fig:dataset-synthetic-examples}
\end{figure}

Pembagian dataset (protokol akademis):

Pengambilan data mengikuti protokol akademis pembagian tiga arah untuk evaluasi yang tidak bias:

\begin{itemize}[nosep]
    \item Pelatihan: 70\% (3.317 citra) --- untuk pembelajaran generator dan diskriminator
    \item Validasi: 15\% (710 citra) --- untuk penyetelan \textit{hyperparameter} dan penghentian dini
    \item Uji: 15\% (712 citra) --- set terpisah untuk evaluasi akhir
\end{itemize}

Dataset diacak dengan seed deterministik (\texttt{seed=42}) sebelum pembagian untuk reproduktibilitas. Jalur degradasi diterapkan dengan randomisasi parameter yang terkontrol, memastikan model belajar pola degradasi umum dan meningkatkan kemampuan generalisasi pada dokumen historis yang tidak terlihat sebelumnya.

Spesifikasi komputasi dan kinerja:

Eksperimen dilakukan pada workstation yang telah dijelaskan pada Tabel~\ref{tab:hardware-specs-actual}. Pelatihan memerlukan sekitar 48 GPU-hours untuk 50 epoch. Inference pada citra baris tunggal ($128\times1024$) memerlukan 50--60 ms, dengan \textit{throughput batch processing} mencapai $\approx$17 citra/detik (batch=2, 30--35 ms per citra) menggunakan 6--8 GB VRAM.

\subsubsection{Dataset Dokumen Historis Nyata (ANRI)}

Untuk evaluasi generalisasi pada degradasi historis nyata, subset dokumen representatif dari Arsip Nasional Republik Indonesia (ANRI) digunakan sebagai set evaluasi kualitatif.

Pemilihan dokumen uji:

Dari koleksi arsip ANRI berupa dokumen administratif VOC dan Hindia Belanda (abad ke-16 hingga ke-18, terutama era 1650--1800), sebanyak 15 dokumen halaman penuh dipilih secara manual untuk evaluasi generalisasi. Pemilihan dilakukan berdasarkan kriteria:

\begin{itemize}[nosep]
    \item Representasi variasi degradasi: Mencakup spektrum tingkat degradasi dari ringan hingga parah untuk menguji \textit{robustness} model
    \item Diversitas jenis degradasi: Penuaan kertas, noda organik, foxing, tinta luntur, bleed-through, kerusakan fisik
    \item Variasi karakteristik paleografi: Berbagai gaya tulisan tangan, tingkat kontras, dan kompleksitas ligatur
\end{itemize}

Karakteristik dataset uji (n=15):

Perbedaan dengan dataset semi-sintetis: Dataset ANRI nyata ini berbeda fundamental dari dataset semi-sintetis untuk evaluasi kuantitatif. Dokumen-dokumen ini mengalami degradasi historis autentik selama berabad-abad penyimpanan arsip, menampilkan pola kerusakan kompleks yang tidak dapat sepenuhnya direplikasi melalui augmentasi komputasional:

\begin{itemize}[nosep]
    \item Degradasi temporal: Penuaan kertas alami dengan penguningan heterogen dan kerapuhan lokal
    \item Artefak organik: Noda biologis, \textit{foxing spots} dengan pola distribusi alami, degradasi tinta \textit{non-uniform}
    \item Kerusakan struktural: Robekan fisik, lipatan permanen, kehilangan material pada tepi dokumen
    \item Variasi tulisan: Tulisan tangan \textit{kursif} paleografi Belanda dengan ligatur kompleks, singkatan periode, dan inkonsistensi gaya penulisan
\end{itemize}

Dokumen ANRI nyata ini tidak memiliki \textit{ground truth} visual berpasangan (citra bersih) karena dokumen historis autentik tidak dapat dipulihkan ke kondisi asli. Oleh karena itu, evaluasi pada dataset nyata bersifat kualitatif melalui inspeksi visual dan validasi ahli paleografi, melengkapi evaluasi kuantitatif pada dataset semi-sintetis. Pendekatan evaluasi dual ini mengikuti praktik standar dalam riset restorasi dokumen historis: dataset semi-sintetis realistis untuk pengukuran kuantitatif presisi, dan dataset nyata dengan degradasi alami untuk validasi generalisasi kualitatif pada kondisi dunia nyata.

\subsection{Implementasi Pipeline Data dan Training Framework}
\label{subsec:implementasi-pipeline}

\subsubsection{Pipeline Pemrosesan Data}

Implementasi framework GAN-HTR memerlukan dua pipeline data terintegrasi yang bekerja secara berurutan: pipeline persiapan dataset HTR menghasilkan citra bersih berpasangan dengan transkripsi \textit{ground truth}, kemudian pipeline degradasi sintetis mentransformasikannya menjadi \textit{dataset triplet} (degraded-clean-text) yang digunakan untuk training. Kedua pipeline dirancang dengan mekanisme validasi kualitas ketat untuk memastikan realisme data dan keselarasan sempurna antara citra terdegradasi dengan ground truth bersihnya.

Pipeline 1: Persiapan Dataset HTR Recognizer. Pipeline ini memproses dokumen historis bersih (halaman penuh) dan file transkripsi XML menjadi dataset TFRecord terstruktur untuk training recognizer. Proses dimulai dengan segmentasi halaman ke baris teks individual menggunakan algoritma deteksi \textit{bounding box} berbasis proyeksi vertikal histogram piksel, menghasilkan koordinat $(x, y, w, h)$ untuk setiap baris yang kemudian diekstrak sebagai \textit{image patch} terpisah. Setiap patch menjalani pra-pemrosesan standar: normalisasi intensitas ke rentang $[0,1]$, \textit{resize bilinear} ke dimensi tetap $128 \times 1024$ piksel (tinggi $\times$ lebar) untuk konsistensi input model, dan binarisasi adaptif Otsu untuk meningkatkan kontras goresan-latar. Validasi kesesuaian memastikan setiap \textit{image patch} memiliki transkripsi teks yang valid (panjang $>$ 3 karakter, tidak mengandung karakter \textit{out-of-vocabulary}). Dataset final di-split dengan rasio 80/10/10 untuk training ($n$=3{,}550 baris), validasi ($n$=710), dan uji ($n$=712) menggunakan \textit{stratified sampling} berdasarkan panjang teks untuk mempertahankan distribusi kompleksitas antar split.

Pipeline 2: Degradasi Sintetis untuk GAN Training. Pipeline ini mentransformasi citra bersih dari Pipeline 1 menjadi \textit{dataset triplet} $(I_{\text{deg}}, I_{\text{clean}}, \text{transcription})$ dengan mengaplikasikan lima jenis degradasi sintetis yang memodelkan kerusakan dokumen historis autentik. Degradasi diterapkan secara independen dengan probabilitas tertentu dan dapat \textit{overlap} untuk menghasilkan kombinasi realistis: (1) bleed-through (30\% probabilitas) mensimulasikan tembus tinta dari halaman balik menggunakan \textit{alpha blending} dengan intensitas acak $\alpha \in [0.3, 0.7]$, (2) \textit{fading} (40\%) mengurangi kontras goresan dengan faktor $\gamma \in [0.5, 0.8]$ untuk mensimulasikan degradasi tinta, (3) \textit{stains} (25\%) menambahkan tekstur noda dengan model Perlin noise dan \textit{Gaussian blur}, (4) \textit{Gaussian blur} (20\%) dengan kernel $\sigma \in [0.5, 2.0]$ mensimulasikan ketidakfokusan optik, dan (5) \textit{salt \& pepper noise} (10\%) dengan densitas $d \in [0.001, 0.01]$ mensimulasikan piksel rusak. Setiap citra terdegradasi divalidasi dengan threshold PSNR $> 20$ dB relatif terhadap \textit{ground truth} untuk memastikan degradasi tidak ekstrem yang dapat merusak struktur karakter fundamental. Pipeline menghasilkan total 3{,}550 \textit{triplet} training, 710 validasi, dan 712 uji dengan keselarasan piksel sempurna antara versi terdegradasi dan bersih---krusial untuk kalkulasi metrik PSNR dan SSIM yang akurat.

\subsection{Strategi dan Protokol Training}
\label{subsec:detail-implementasi-training}

\subsubsection{Curriculum Learning dan Protokol Optimisasi}

Pelatihan menggunakan pendekatan curriculum learning tiga fase untuk stabilitas optimal dan konvergensi bertahap:

\begin{enumerate}
    \item Fase \textit{Warmup} Visual (epoch 1--10): Pelatihan eksklusif pada kualitas visual tanpa supervisi teks, memungkinkan generator membangun kemampuan rekonstruksi dasar. Fase ini mencegah mode collapse dini yang dapat terjadi jika loss CTC diperkenalkan terlalu awal. Pada fase ini, hanya komponen loss visual yang diaktifkan: $\mathcal{L}_{\text{pixel}}$, $\mathcal{L}_{\text{adv}}$, dan $\mathcal{L}_{\text{perc}}$ dengan bobot penuh sesuai Tabel~\ref{table:optimal_weights}, sementara $\lambda_{\text{ctc}}=0$ untuk stabilitas awal.
    
    \item Fase Peningkatan Bertahap CTC (epoch 11--30): Peningkatan bertahap bobot CTC dari nol ke bobot penuh secara linear. Integrasi halus ini memungkinkan generator beradaptasi dengan gradien HTR tanpa mengganggu kestabilan pelatihan adversarial yang sudah terbentuk. Bobot CTC ditingkatkan secara linear: $\lambda_{\text{ctc}}(e) = 0.15 \times \frac{e - 10}{20}$ untuk epoch $e \in [11, 30]$, mencapai nilai maksimal 0.15 pada epoch 30.
    
    \item Fase Pelatihan Penuh (epoch 31--50): Pelatihan dengan semua komponen loss aktif pada bobot optimal (Tabel~\ref{table:optimal_weights}), disertai penyeimbangan adaptif untuk menjaga keseimbangan dinamis antara loss visual dan tekstual. Mekanisme penyeimbangan memantau rasio kontribusi setiap komponen dan melakukan penyesuaian minor ($\pm$5\%) jika terjadi dominasi berlebihan dari satu komponen.
\end{enumerate}

\paragraph{Prosedur Pelatihan \textit{Adversarial} Bergantian}

Prosedur pelatihan adversarial bergantian dengan supervisi dual-modal diterapkan pada setiap iterasi:

\begin{enumerate}
    \item Langkah Diskriminator: Perbarui $D$ untuk memaksimalkan kemampuan membedakan pasangan nyata $(I_{\text{gt}}, \text{teks}_{\text{gt}})$ dari pasangan palsu $(I_{\text{gen}}, \text{teks}_{\text{pred}})$ dengan supervisi simultan pada domain visual dan tekstual. Diskriminator dilatih dengan loss \textit{binary cross-entropy} yang mengoptimalkan deteksi autentisitas pada kedua modalitas secara bersamaan.
    
    \item Langkah Generator: Perbarui $G$ untuk meminimalkan $\mathcal{L}_{\text{total}}$ (kombinasi lima komponen loss, Persamaan~\ref{eq:IV.12}), menghasilkan citra yang secara simultan menipu diskriminator dan mengoptimalkan keterbacaan HTR. Generator menerima gradien dari diskriminator (untuk realisme visual), recognizer beku (untuk keterbacaan teks), dan loss rekonstruksi langsung.
\end{enumerate}

Pembaruan dilakukan dengan rasio seimbang 1:1, memastikan stabilitas pelatihan adversarial. Setiap iterasi pelatihan melakukan satu pembaruan diskriminator diikuti satu pembaruan generator untuk mencegah ketidakseimbangan yang dapat menyebabkan mode collapse atau dominasi salah satu komponen.

\paragraph{Konfigurasi Pelatihan}

Konfigurasi pelatihan didasarkan pada penyetelan empiris untuk stabilitas GAN dan konvergensi optimal:

\begin{itemize}[nosep]
    \item Optimisasi: Generator menggunakan algoritma Adam dengan parameter momentum disesuaikan ($\beta_1=0.5$, $\beta_2=0.999$) untuk stabilitas GAN, sementara diskriminator menggunakan SGD dengan momentum~$=0.9$ untuk mencegah mode collapse dan memastikan konvergensi stabil. Kedua model menerapkan gradient clipping dengan norma global~$=1.0$ untuk mencegah ledakan gradien.
    
    \item Laju Pembelajaran: Kedua model menggunakan laju pembelajaran awal yang sama ($\alpha_{\text{init}}=2\times10^{-4}$) dengan penjadwalan \textit{cosine annealing} untuk konvergensi halus tanpa penurunan tajam. Laju pembelajaran menurun secara halus mengikuti fungsi \textit{cosine} dari nilai awal ke nilai minimum ($\alpha_{\text{min}}=1\times10^{-6}$) selama pelatihan.
    
    \item Ukuran \textit{Batch} dan Dimensi: Ukuran batch kecil (2 sampel per GPU) dengan dimensi citra lebar ($1024 \times 128 \times 1$ piksel) yang mempertahankan detail goresan halus pada teks baris tunggal paleografi. Dimensi lebar diperlukan untuk menangkap konteks sekuensial teks baris lengkap.
    
    \item Penghentian Dini: Diterapkan dengan mekanisme \textit{patience} 25 epoch yang mempertimbangkan fase \textit{curriculum}, memulihkan bobot terbaik berdasarkan metrik gabungan PSNR dan CER dengan pembobotan yang memprioritaskan kualitas visual. Kriteria: $\text{Score} = \text{PSNR} - 0.5 \times \text{CER}$ untuk keseimbangan optimal.
\end{itemize}

\begin{table}[H]
\caption[Protokol curriculum learning tiga fase dengan strategi pelatihan]{Protokol curriculum learning tiga fase.\\
Strategi pelatihan bertahap dengan penjadwalan bobot CTC loss ($\lambda_{\text{CTC}}$) untuk stabilitas konvergensi.\\
Fase Warmup Visual (epoch 1--10, $\lambda_{\text{CTC}}=0$) membangun struktur visual dasar,\\
Fase Peningkatan Bertahap (epoch 11--30, $\lambda_{\text{CTC}}: 0 \to 0.15$) mengintegrasikan gradien HTR secara halus,\\
Fase Pelatihan Penuh (epoch 31--50, $\lambda_{\text{CTC}}=0.15$) mengoptimalkan semua komponen secara simultan.}
\label{tab:curriculum-learning}
\small
\centering
\begin{tabular}{|l|p{2cm}|p{2.5cm}|l|}
\hline
\textbf{Fase} & \textbf{Epoch} & \textbf{$\lambda_{CTC}$} & \textbf{Fokus} \\
\hline
\textit{Warmup} Visual & 1--10 & 0.0 & Struktur visual dasar, stabilitas adversarial \\
Peningkatan Bertahap & 11--30 & $0 \to 0.15$ & Adaptasi gradien HTR secara halus \\
Pelatihan Penuh & 31--50 & 0.15 & Optimasi penuh semua komponen \\
\hline
\end{tabular}
\end{table}

\subsubsection{Analisis Stabilitas dan Regularisasi}

Gambar~\ref{fig:convergence_analysis} menunjukkan trajektori konvergensi pelatihan selama 50 epoch, memvalidasi stabilitas dan efektivitas strategi pelatihan yang diusulkan.

\begin{figure}[H]
\centering
\includegraphics[width=0.95\textwidth]{/home/lambda_one/tesis/GAN-HTR-ORI/docRestoration/Paper/data_dukung/fig_convergence_analysis_combined.pdf}
\caption[Trajektori konvergensi pelatihan selama 50 epoch]{Trajektori konvergensi pelatihan selama 50 epoch. (a) \textit{Loss generator} menunjukkan penurunan stabil tanpa mode collapse, (b) PSNR meningkat mencapai target yang ditetapkan, (c) SSIM mencapai nilai tinggi menunjukkan pelestarian struktur. Model terbaik (ditandai garis vertikal putus-putus dan bintang merah) dipilih berdasarkan kriteria gabungan PSNR-CER. Kurva konvergensi halus memvalidasi efektivitas strategi curriculum learning yang diusulkan dengan tiga fase: \textit{Warmup} Visual (epoch 1--10), Peningkatan Bertahap CTC (epoch 11--30), dan Pelatihan Penuh (epoch 31--50).}
\label{fig:convergence_analysis}
\end{figure}

Observasi konvergensi:

Analisis trajektori konvergensi mengungkapkan pola karakteristik yang mengonfirmasi stabilitas pelatihan:

\begin{itemize}[nosep]
    \item \textit{Loss Generator}: Penurunan monotonik halus pada fase awal, stabil pada fase akhir tanpa osilasi, mengindikasikan pelatihan adversarial yang stabil tanpa mode collapse. Tidak terdeteksi lonjakan loss yang mengindikasikan divergensi atau ketidakstabilan gradien.
    
    \item Trajektori PSNR: Peningkatan cepat pada Fase \textit{Warmup} Visual (epoch 1--10), diikuti peningkatan stabil pada Fase Peningkatan Bertahap CTC (epoch 11--30), dan \textit{plateau} pada fase akhir yang memicu pertimbangan penghentian dini. Target PSNR $>$30~dB tercapai pada epoch 38, mengonfirmasi efektivitas strategi \textit{curriculum}.
    
    \item Progresi SSIM: Tren naik konsisten mencapai nilai tinggi (0.987) pada epoch terbaik, menunjukkan pelestarian kesamaan struktural bersama akurasi tingkat piksel. Konvergensi SSIM lebih cepat dibanding PSNR, mencapai target $>$0.95 pada epoch 22.
    
    \item Pemilihan Model Terbaik: Dipilih berdasarkan optimisasi metrik gabungan PSNR-CER pada epoch 44 (bukan penghentian arbitrer pada epoch terakhir), mengikuti protokol \textit{early stopping} dengan \textit{patience} 25 epoch. Model ini menyeimbangkan kualitas visual dan keterbacaan teks secara optimal.
    
    \item Keandalan Statistik: Interval kepercayaan yang sempit (ditunjukkan oleh area bayangan pada grafik) memvalidasi kinerja robust pada validation set, mengindikasikan konsistensi prediksi model tanpa fluktuasi berlebihan antar batch.
\end{itemize}

Teknik regularisasi:

Untuk mencegah overfitting dan meningkatkan generalisasi model, teknik regularisasi berikut diterapkan:

\begin{itemize}[nosep]
    \item \textit{Batch Normalization}: Diterapkan di semua lapisan konvolusional dengan momentum~$=0.9$ dan $\epsilon=10^{-5}$ untuk stabilitas pelatihan dan mempercepat konvergensi. Normalisasi memastikan distribusi aktivasi tetap stabil selama pelatihan, mengurangi \textit{internal covariate shift}.
    
    \item \textit{Label Smoothing} Asimetris Satu Sisi: Hanya diterapkan pada label nyata ($\epsilon_{\text{smooth}}=0.1$, mengubah label 1.0 menjadi 0.9) untuk mengurangi kepercayaan berlebihan diskriminator sambil mempertahankan sinyal pembelajaran yang memadai untuk sampel palsu (label 0.0 tidak diubah). Pendekatan asimetris ini mencegah diskriminator menjadi terlalu yakin yang dapat menghambat pembelajaran generator.
    
    \item Koneksi Residual: Diterapkan di diskriminator untuk aliran gradien yang lebih baik dan pelatihan arsitektur dalam yang stabil. Koneksi \textit{skip} memungkinkan gradien mengalir langsung ke lapisan awal, mengatasi masalah \textit{vanishing gradient} pada arsitektur \textit{deep}.
    
    \item \textit{Gradient Clipping}: Diterapkan pada kedua model dengan norma global~$=1.0$ untuk mencegah ledakan gradien, khususnya dari loss CTC yang di-\textit{clip} pada nilai maksimum~$=400.0$ untuk stabilitas numerik. \textit{Clipping} CTC mencegah gradien ekstrem yang dapat muncul dari kesalahan prediksi besar pada teks yang sangat terdegradasi.
    
    \item \textit{Dropout}: Diterapkan pada \textit{classification head} diskriminator (\textit{dropout rate}~$=0.1$) dan \textit{embedding layer} teks (\textit{dropout rate}~$=0.2$) untuk regularisasi tambahan. \textit{Dropout} mencegah ko-adaptasi berlebihan antar neuron, meningkatkan kemampuan generalisasi model.
\end{itemize}

Detail implementasi teknis:

Spesifikasi parameter pelatihan lengkap untuk reproduktifitas:

\begin{itemize}[nosep]
    \item Optimizer Diskriminator: SGD dengan momentum~$=0.9$, laju pembelajaran awal $\alpha_{\text{init}}=2\times10^{-4}$
    \item \textit{Optimizer Generator}: Adam dengan $\beta_1=0.5$, $\beta_2=0.999$, laju pembelajaran awal $\alpha_{\text{init}}=2\times10^{-4}$
    \item Ukuran \textit{Batch}: 2 sampel per GPU (total 4 dengan 2 GPU)
    \item Dimensi Citra: $1024 \times 128 \times 1$ piksel (\textit{width} $\times$ \textit{height} $\times$ \textit{channel})
    \item Presisi Numerik: \textit{Pure} FP32 (tanpa mixed precision) untuk stabilitas CTC loss
    \item \textit{Early Stopping}: \textit{Patience}~$=25$ epoch, \textit{min delta}~$=0.05$ pada metrik gabungan PSNR-CER
    \item \textit{Gradient Clipping}: Norma global~$=1.0$ untuk kedua model
    \item CTC loss Clipping: Nilai maksimum~$=400.0$ untuk stabilitas numerik
\end{itemize}

Catatan penghentian dini:

Model terbaik dipilih berdasarkan metrik gabungan PSNR-CER pada epoch 44 yang menghasilkan kinerja optimal. Pelatihan dilanjutkan hingga epoch 50 untuk memverifikasi tidak ada peningkatan signifikan. Mekanisme penghentian dini tidak terpicu karena fluktuasi kecil dalam metrik validasi masih berada dalam ambang batas yang ditentukan (\textit{min delta} 0.05), mengonfirmasi model telah mencapai konvergensi stabil. Pemilihan epoch 44 (bukan 50) memastikan model tidak mengalami overfitting pada fase pelatihan akhir.

\paragraph{Implementasi dan Validasi Mode Diskriminator}

Diskriminator Dual-Modal mendukung dua mode input teks: (1) Ground Truth Mode menggunakan teks ground truth, memberikan stabilitas superior (konvergensi 100\% vs 85\%, target PSNR tercapai epoch 38 vs 45, deviasi standar validation loss 35\% lebih rendah), (2) Predicted Mode menggunakan prediksi frozen recognizer, lebih realistis untuk deployment namun menantang dengan CER 33.72\% (perlu durasi $\sim$1.15x epoch, pemantauan stabilitas cermat). Perbandingan kuantitatif detail dilaporkan pada Bab V.4.3 (Ablasi \textit{Diskriminator Mode}).

Implementasi final menggunakan Ground Truth Mode berdasarkan validasi eksperimen (5 \textit{independent runs}): performa sebanding (perbedaan PSNR $<$0.3 dB, $p>0.05$) namun Predicted Mode konvergensi 15\% lebih rendah dengan degradasi CER validasi +1.2\%. Keputusan ini memprioritaskan stabilitas pelatihan untuk eksperimen ablasi sistematis (Bab V) sambil mempertahankan opsi Predicted Mode untuk deployment produksi.

\paragraph{Kebijakan Presisi Numerik}

Kebijakan presisi pure FP32 diimplementasikan secara eksplisit untuk memastikan stabilitas numerik dalam pelatihan, terutama untuk kalkulasi \textit{CTC loss} yang memerlukan presisi penuh.

Alasan kebijakan FP32:

Pemilihan pure FP32 (tanpa \textit{mixed precision FP16}) didasarkan pada analisis stabilitas numerik:

\begin{enumerate}[nosep]
\item Presisi Numerik CTC loss: Algoritma CTC melakukan komputasi \textit{log-probability} pada \textit{sequence alignment} yang sensitif terhadap \textit{underflow}. Eksperimen awal dengan \textit{mixed precision FP16} (diuji pada 3 \textit{runs}) menunjukkan \textit{gradient underflow} pada epoch 15--20, menyebabkan divergensi pelatihan pada 1 dari 3 \textit{runs} (tingkat kegagalan 33\%).

\item Konsistensi Multi-Component \textit{Loss}: Dengan 5 komponen loss berbeda (adversarial, L1, \textit{perceptual}, CTC, \textit{rec-feat}), presisi konsisten diperlukan untuk penyeimbangan yang stabil. FP16 dapat menyebabkan inkonsistensi numerik antar komponen karena perbedaan magnitudo (\textit{adversarial loss} $\sim$0.1--1.0, \textit{L1 loss} $\sim$10--50).

\item Stabilitas \textit{Gradient Flow}: \textit{Pure FP32} memastikan \textit{backpropagation} melalui frozen recognizer (27.86M parameter) tidak mengalami kehilangan presisi yang dapat merusak gradien generator.
\end{enumerate}

Kebijakan pure FP32 dipilih untuk stabilitas numerik \textit{CTC loss} dan reproduktifitas, dengan trade-off waktu pelatihan $\sim$28\% lebih lambat dibanding \textit{mixed precision FP16}. Validasi empiris analisis biaya-manfaat dilaporkan pada Bab V Subbab 5.2 (Ablasi Presisi Numerik).

\subsection{Metode Baseline dan Strategi Evaluasi}
\label{subsec:metode-baseline-strategi}

Evaluasi metode yang diusulkan dilakukan melalui dua pendekatan komplementer: (1) perbandingan kuantitatif dengan metode klasik yang dapat direplikasi pada kondisi eksperimental identik, dan (2) studi ablasi sistematis untuk memvalidasi kontribusi setiap komponen arsitektural.

\subsubsection{Metode Klasik untuk Perbandingan Kuantitatif}

Metode klasik dipilih sebagai baseline kuantitatif karena dapat diimplementasi ulang pada dataset identik tanpa memerlukan pelatihan ulang. Tiga metode baseline digunakan dengan detail implementasi sebagai berikut:

Baseline 1: Tanpa Perbaikan (Degraded Input)

Evaluasi HTR langsung pada citra terdegradasi tanpa pra-pemrosesan untuk mengukur dampak degradasi terhadap kinerja recognizer:

\begin{itemize}[nosep]
    \item Pra-pemrosesan: Normalisasi piksel ke rentang [0, 1] (identik dengan pipeline pelatihan recognizer)
    \item Input: Citra grayscale $128 \times 1024$ piksel dari test set ($n$=712)
    \item HTR Engine: \textit{Frozen recognizer} yang sama digunakan pada framework yang diusulkan
    \item Tujuan: Mengukur \textit{upper bound} degradasi (CER/WER terburuk) sebagai referensi perbaikan
\end{itemize}

Baseline 2: Otsu Thresholding

Binarisasi global berbasis histogram untuk pembersihan dokumen dengan implementasi OpenCV:

\begin{itemize}[nosep]
    \item Library: OpenCV 4.8.1
    \item Method: Otsu's automatic thresholding
    \item Pra-pemrosesan:
    \begin{enumerate}[nosep]
        \item Gaussian blur dengan kernel $5 \times 5$ untuk reduksi derau
        \item Konversi ke 8-bit unsigned integer [0, 255]
        \item Otsu thresholding (threshold otomatis dari histogram)
    \end{enumerate}
    \item Pasca-pemrosesan: Normalisasi kembali ke [0, 1] untuk input HTR
    \item Parameter: Threshold ditentukan otomatis oleh algoritma Otsu, tidak ada \textit{tuning} manual
\end{itemize}

Implementasi standar tanpa optimasi parameter memastikan perbandingan yang adil dengan metode \textit{off-the-shelf}.

Baseline 3: Sauvola Adaptive Thresholding

Binarisasi adaptif lokal untuk dokumen dengan pencahayaan \textit{non-uniform} menggunakan implementasi scikit-image:

\begin{itemize}[nosep]
    \item Library: scikit-image 0.21.0
    \item Parameters: Window size $w = 51$ piksel (dioptimasi via grid search), parameter standar $k = 0.2$
    \item Pra-pemrosesan: Gaussian blur kernel $3 \times 3$ (lebih ringan dari Otsu)
    \item Optimization: Parameter $w$ dioptimalkan melalui pencarian kisi $w \in \{25, 51, 75, 101\}$ pada \textit{subset validation} ($n$=100), dipilih $w=51$ berdasarkan CER terendah
\end{itemize}

\paragraph{Justifikasi Pemilihan \textit{Baseline} Klasik}

Metode klasik dipilih sebagai baseline kuantitatif karena: (1) dapat diimplementasi ulang pada dataset identik tanpa pelatihan ulang, (2) parameter yang dapat direproduksi dengan library standar, (3) waktu eksekusi cepat untuk evaluasi batch ($<$1 detik per citra), (4) mewakili pendekatan tradisional yang umum digunakan praktisi arsip. Perbandingan dengan metode \textit{deep learning state-of-the-art} (DE-GAN, DocEnTr) tidak dilakukan karena perbedaan dataset pelatihan dan ketiadaan implementasi resmi yang dapat direproduksi.

\subsubsection{Metode Deep Learning Terkini (Konteks Penelitian)}

Metode berbasis \textit{deep learning} yang dibahas dalam tinjauan literatur (Bab II) berfungsi sebagai konteks penelitian untuk memposisikan kontribusi yang diusulkan:

\begin{itemize}[nosep]
    \item U-Net: Arsitektur \textit{encoder-decoder} dengan skip connections untuk restorasi citra
    \item DE-GAN: GAN bersyarat untuk peningkatan dokumen dengan fokus kualitas visual
    \item HTR-GAN: GAN dengan pengenal yang dilatih bersama untuk restorasi berorientasi HTR
    \item DocEnTr: Arsitektur berbasis Transformer dengan mekanisme atensi mandiri untuk peningkatan dokumen
\end{itemize}

Perbandingan kuantitatif langsung dengan metode \textit{deep learning} terkini tidak dilakukan karena: (1) metode tersebut dilatih pada dataset berbeda dengan karakteristik degradasi yang tidak identik, sehingga perbandingan langsung tidak seimbang, (2) implementasi resmi yang dapat direproduksi tidak tersedia untuk sebagian besar metode, dan (3) perbedaan arsitektur, protokol pelatihan, dan fungsi loss membuat isolasi kontribusi spesifik sulit dilakukan. Pendekatan ini mengikuti praktik evaluasi yang umum dalam riset restorasi dokumen di mana perbandingan dilakukan hanya pada kondisi eksperimental yang dapat direplikasi secara identik.

\subsubsection{Studi Ablasi sebagai Validasi Komponen}

Untuk memvalidasi kontribusi setiap komponen arsitektural secara sistematis, studi ablasi inkremental dilakukan dengan menambahkan komponen loss secara bertahap (detail evaluasi pada Bab V):

\begin{enumerate}[nosep]
    \item Eksperimen 01 (Hanya Piksel): Setara dengan U-Net standar untuk restorasi citra
    \item Eksperimen 02 (+\textit{Adversarial}): Setara dengan GAN standar tanpa komponen HTR
    \item Eksperimen 03 (+Perseptual): Peningkatan dengan loss perseptual berbasis VGG
    \item Eksperimen 04 (+CTC): Konfigurasi optimal dengan integrasi HTR
    \item Eksperimen 05 (+Fitur Pengenal): Konfigurasi lengkap untuk validasi redundansi komponen
\end{enumerate}

Pendekatan ablasi ini memberikan pemahaman komprehensif tentang kontribusi relatif setiap komponen dan memvalidasi keputusan desain arsitektural secara empiris.

\subsection{Konfigurasi Metrik Evaluasi dan Protokol Eksperimen}
\label{subsec:konfigurasi-metrik-evaluasi}

Evaluasi framework GAN-HTR yang diusulkan memerlukan metrik multi-objektif untuk mengukur kualitas visual dan keterbacaan teks secara simultan, serta protokol eksperimen yang terstruktur untuk memastikan reproduktibilitas. Tabel~\ref{tab:metric-component-mapping} memetakan metrik evaluasi dengan komponen arsitektur yang divalidasi dan kebutuhan sistem yang telah ditetapkan pada Bab III.

\begin{table}[H]
\caption[Pemetaan metrik evaluasi dengan komponen arsitektur dan kebutuhan sistem]{Pemetaan metrik evaluasi dengan komponen arsitektur dan kebutuhan sistem. Tabel menunjukkan korespondensi antara empat metrik utama (PSNR, SSIM, CER, WER) dengan aspek kualitas yang diukur, komponen arsitektur terkait (generator, frozen recognizer), dan kebutuhan fungsional/non-fungsional sistem yang dipenuhi. Pemetaan ini memvalidasi desain loss function yang selaras dengan tujuan evaluasi.}
\label{tab:metric-component-mapping}
\small
\centering
\begin{tabular}{|l|l|l|l|}
\hline
\textbf{Metrik} & \textbf{Aspek Diukur} & \textbf{Komponen Terkait} & \textbf{Kebutuhan} \\
\hline
PSNR & Akurasi piksel & Generator ($\mathcal{L}_{\text{pixel}}$) & NFR-1 \\
SSIM & Struktur perseptual & Generator ($\mathcal{L}_{\text{perc}}$) & NFR-1 \\
CER & Keterbacaan karakter & Frozen Recognizer + $\mathcal{L}_{\text{CTC}}$ & FR-3, NFR-2 \\
WER & Keterbacaan kata & Frozen Recognizer + $\mathcal{L}_{\text{CTC}}$ & FR-3, NFR-2 \\
\hline
\end{tabular}
\end{table}

\subsubsection{Metrik Kualitas Visual}

Detail implementasi teknis untuk metrik kualitas visual dengan spesifikasi library, versi, dan parameter konkret.

PSNR (\textit{Peak Signal-to-Noise Ratio}):
\begin{itemize}[nosep]
    \item Library: scikit-image 0.21.0
    \item Formula: 
    \begin{equation}
    \text{PSNR} = 10 \log_{10}\left(\frac{\text{MAX}^2}{\text{MSE}}\right)
    \label{eq:IV.17}
    \end{equation}
    dengan $\text{MAX}=1.0$ untuk citra normalisasi $[0,1]$, dan $\text{MSE}$ adalah \textit{Mean Squared Error}
    \item Pra-pemrosesan: Citra ternormalisasi ke format float32 $[0,1]$
    \item Target: PSNR $>$ 30~dB (kualitas restorasi tinggi, sebagaimana ditetapkan Bab III)
\end{itemize}

SSIM (\textit{Structural Similarity Index}):
\begin{itemize}[nosep]
    \item Library: scikit-image 0.21.0
    \item Konfigurasi: Standar untuk citra grayscale ternormalisasi $[0,1]$ dengan bobot Gaussian
    \item Formula: 
    \begin{equation}
    \text{SSIM}(I_1, I_2) = \frac{(2\mu_1\mu_2 + C_1)(2\sigma_{12} + C_2)}{(\mu_1^2 + \mu_2^2 + C_1)(\sigma_1^2 + \sigma_2^2 + C_2)}
    \label{eq:IV.18}
    \end{equation}
    dengan $I_1$, $I_2$ adalah citra yang dibandingkan; $\mu_1$, $\mu_2$ adalah rata-rata lokal; $\sigma_1$, $\sigma_2$ adalah variansi lokal; $\sigma_{12}$ adalah kovariansi; $C_1$, $C_2$ adalah konstanta stabilisasi
    \item Target: SSIM $> 0.95$ (kesamaan struktural tinggi, Bab III)
\end{itemize}

\subsubsection{Metrik Keterbacaan Teks}

CER (\textit{Character Error Rate}):
\begin{itemize}[nosep]
    \item Library: editdistance 0.6.2
    \item Algoritma: \textit{Levenshtein distance}
    \item Formula: 
    \begin{equation}
    \text{CER} = \frac{\text{EditDistance}(t_{\text{gt}}, t_{\text{pred}})}{|t_{\text{gt}}|} \times 100\%
    \label{eq:IV.19}
    \end{equation}
    dengan $t_{\text{gt}}$ adalah teks ground truth, $t_{\text{pred}}$ adalah teks prediksi, dan $|t_{\text{gt}}|$ adalah panjang karakter \textit{ground truth}
    \item Pra-pemrosesan: CTC greedy decoding dengan normalisasi lowercase dan penghapusan \textit{whitespace} berlebih
    \item Target: Penurunan signifikan dari baseline terdegradasi (CER $\sim$83\%)
\end{itemize}
WER (\textit{Word Error Rate}):
\begin{itemize}[nosep]
    \item Library: editdistance 0.6.2
    \item Formula: 
    \begin{equation}
    \text{WER} = \frac{\text{EditDistance}(w_{\text{gt}}, w_{\text{pred}})}{|w_{\text{gt}}|} \times 100\%
    \label{eq:IV.20}
    \end{equation}
    dengan $w_{\text{gt}}$ adalah urutan kata \textit{ground truth}, $w_{\text{pred}}$ adalah urutan kata prediksi, dan $|w_{\text{gt}}|$ adalah jumlah kata \textit{ground truth}
    \item Pra-pemrosesan: Identik dengan CER dengan tokenisasi tingkat kata
\end{itemize}

\subsubsection{Klarifikasi Baseline CER Recognizer}

Tabel~\ref{table:cer-summary-baseline} merangkum tiga nilai baseline CER yang digunakan dalam penelitian ini: CER$_{\text{pretrain}}$ (33.72\%) pada dataset ANRI awal untuk validasi generalisasi paleografi Belanda, CER$_{\text{val-GT}}$ (26.57\%) pada validation set penelitian, dan CER$_{\text{test-GT}}$ (34.1\%) pada test set penelitian yang menjadi referensi utama evaluasi. Perbedaan CER$_{\text{val-GT}}$ vs CER$_{\text{test-GT}}$ ($\Delta$=7.53 poin persentase) mencerminkan variasi kompleksitas paleografi antar dataset (ligatur lebih kompleks, variasi gaya tulisan lebih besar pada test set), bukan overfitting---dikonfirmasi oleh konsistensi \textit{gap} pada hasil restorasi: CER$_{\text{val-restored}}$ = 27.11\% vs CER$_{\text{test-restored}}$ = 34.9\% ($\Delta$=7.79\%), divalidasi pada Bab V.3.2.

\begin{table}[H]
\caption[Ringkasan metrik CER pada berbagai dataset dan kondisi.]{Ringkasan metrik CER pada berbagai dataset dan kondisi. Perbandingan performa frozen recognizer pada citra bersih (baseline CER 33.72\% pada ANRI validasi awal), input terdegradasi tanpa restorasi (CER 83.4\%), dan hasil restorasi metode usulan (CER 34.9\% pada test set). Improvement relatif mencapai 58.2\% pada test set, mengonfirmasi efektivitas framework dalam memulihkan keterbacaan teks dari dokumen terdegradasi berat.}
\label{table:cer-summary-baseline}
\small
\centering
\begin{tabular}{|l|p{3.5cm}|p{3.5cm}|p{2cm}|}
\hline
\textbf{Kondisi} & \textbf{Dataset} & \textbf{$n$} & \textbf{CER (\%)} \\
\hline
\multicolumn{4}{|c|}{\textit{\textbf{Baseline Recognizer} pada Citra Bersih}} \\
\hline
CER$_{\text{pretrain}}$ & ANRI validasi awal & 948 & 33.72 \\
CER$_{\text{val-GT}}$ & Validasi GT penelitian & 710 & 26.57 \\
CER$_{\text{test-GT}}$ & Uji GT penelitian & 712 & 34.1 \\
\hline
\multicolumn{4}{|c|}{\textit{\textbf{Input Terdegradasi} (Tanpa Restorasi)}} \\
\hline
CER$_{\text{degraded}}$ & Validasi \& Uji & 1422 & 83.4 \\
\hline
\multicolumn{4}{|c|}{\textit{\textbf{Hasil Restorasi} Metode Usulan}} \\
\hline
CER$_{\text{val-restored}}$ & Validasi & 710 & 27.11 \\
CER$_{\text{test-restored}}$ & Uji & 712 & \textbf{34.9} \\
\hline
\multicolumn{4}{|c|}{\textit{\textbf{Improvement} Relatif}} \\
\hline
Validasi & 83.4\% $\rightarrow$ 27.11\% & - & 67.5\% \\
Uji & 83.4\% $\rightarrow$ 34.9\% & - & \textbf{58.2\%} \\
\hline
\end{tabular}
\end{table}

Tabel~\ref{table:cer-summary-baseline} merangkum semua nilai CER yang dilaporkan dalam penelitian ini untuk transparansi dan reproduktibilitas. Penggunaan CER$_{\text{test-GT}}$ = 34.1\% sebagai referensi utama konsisten dengan protokol evaluasi standar yang menggunakan test set sebagai metrik akhir.

\subsubsection{Protokol Pelaporan Statistik dan Evaluasi}

Semua metrik dilaporkan sebagai rata-rata $\pm$ simpangan baku dengan interval kepercayaan 95\%, dihitung dari set validasi lengkap ($n$=710) atau set uji ($n$=712). Signifikansi statistik untuk perbandingan model diuji menggunakan t-test sampel independen dengan $\alpha = 0.05$. Semua eksperimen menggunakan seed tetap untuk reproduktibilitas penuh.

Protokol Evaluasi Final: Semua hasil yang dilaporkan dalam penelitian ini menggunakan evaluasi konsisten pada checkpoint terbaik (epoch 44, dipilih berdasarkan metrik validasi dengan early stopping patience=25) yang dievaluasi pada set uji held-out ($n$=712, tidak pernah dilihat selama pelatihan). Pendekatan ini mengikuti protokol akademik standar untuk mencegah overfitting dan memastikan hasil yang tidak bias.

\subsubsection{Metrik Diagnostik (Training Monitoring)}

Generator Loss dan Diskriminator Loss:
\begin{itemize}[nosep]
    \item Framework: TensorFlow 2.16.1 dengan \textit{eager execution}
    \item Tracking: TensorBoard + MLflow untuk visualisasi dan \textit{logging}
    \item Logging Frequency: Setiap batch untuk analisis granular, agregasi per epoch untuk tren
    \item Tujuan: Pemantauan stabilitas pelatihan, deteksi mode collapse, validasi konvergensi
    \item Catatan: Metrik diagnostik tidak dilaporkan sebagai metrik evaluasi utama (hanya untuk \textit{monitoring})
\end{itemize}

\subsubsection{Spesifikasi Eksperimen Ablasi}

Detail konfigurasi untuk studi ablasi sistematis yang diuraikan secara konseptual pada Bab III, dengan fokus pada reproduktibilitas dan isolasi kontribusi komponen.

Ablasi Komponen Loss Function:

Protokol Eksperimen:
\begin{itemize}[nosep]
    \item Konfigurasi: Lima eksperimen inkremental (Exp 01--05) dengan penambahan komponen loss bertahap
    \item Dataset Split: Identik untuk semua eksperimen (train 70\%, val 15\%, test 15\%, seed=42)
    \item Durasi: 15 epoch untuk ablasi cepat (vs 50 \textit{epoch production training})
    \item Evaluasi: \textit{Validation set} ($n$=710) untuk iterasi cepat, test set ($n$=712) untuk konfigurasi optimal
    \item Checkpointing: \textit{Best checkpoint} berdasarkan \textit{combined score} PSNR-CER pada validasi
    \item Hyperparameter: Identik untuk semua eksperimen (LR, \textit{batch size}, optimizer) untuk isolasi dampak \textit{loss}
\end{itemize}

Konfigurasi Spesifik Per Eksperimen:
\begin{table}[H]
\caption[Konfigurasi lima eksperimen ablasi komponen loss dengan desain studi]{Konfigurasi lima eksperimen ablasi komponen loss.Desain studi ablasi bertahap untuk mengisolasi kontribusi setiap komponen loss terhadap performa restorasi. Dimulai dari baseline U-Net dengan pixel loss saja (Exp 01), kemudian menambahkan adversarial loss (Exp 02), perceptual loss (Exp 03), CTC loss (Exp 04, konfigurasi optimal), dan recognition feature loss (Exp 05). Protokol ini memungkinkan analisis kontribusi incremental setiap komponen terhadap metrik PSNR, SSIM, dan CER.}
\label{tab:ablasi-loss-config}
\small
\centering
\begin{tabularx}{\textwidth}{|l|X|X|X|X|X|}
\hline
\textbf{Eksperimen} & \textbf{Pixel} & \textbf{Adv} & \textbf{Perc} & \textbf{CTC} & \textbf{RecFeat} \\
\hline
Exp 01 (Baseline U-Net) & $\lambda=50$ & - & - & - & - \\
Exp 02 (+Adversarial) & $\lambda=50$ & $\lambda=3$ & - & - & - \\
Exp 03 (+Perceptual) & $\lambda=50$ & $\lambda=3$ & $\lambda=1$ & - & - \\
Exp 04 (+CTC, Optimal) & $\lambda=50$ & $\lambda=3$ & $\lambda=1$ & $\lambda=0.15$ & - \\
Exp 05 (+RecFeat) & $\lambda=50$ & $\lambda=3$ & $\lambda=1$ & $\lambda=0.15$ & $\lambda=8$ \\
\hline
\end{tabularx}
\end{table}

Ablasi Arsitektur Diskriminator:

Protokol Eksperimen:
\begin{itemize}[nosep]
    \item Konfigurasi: Dua eksperimen dengan generator identik, diskriminator berbeda
    \item Varian 1: Diskriminator CNN \textit{single-modal} (7.2M parameter, baseline)
    \item Varian 2: Diskriminator dual-modal CNN+BiLSTM (17.4M parameter, yang diusulkan)
    \item Durasi: 50 \textit{epoch production training} (bukan ablasi singkat)
    \item Loss Configuration: Identik 4-komponen (Pixel + Adv + Perc + CTC, tanpa RecFeat)
    \item Evaluasi: \textit{Best checkpoint epoch} 44 pada validation set berdasarkan \textit{combined score}
\end{itemize}

Reproduktibilitas dan Determinisme:

Konfigurasi untuk memastikan hasil yang dapat direproduksi:
\begin{itemize}[nosep]
    \item Random Seed: Seed tetap (42) untuk NumPy, TensorFlow, Python
    \item Deterministic Operations: Mode deterministik TensorFlow 2.16+ dan cuDNN
    \item GPU Configuration: Single GPU (GPU:0) untuk konsistensi komputasi
    \item Dependency Management: Poetry untuk version locking semua library
    \item Experiment Tracking: MLflow dengan auto-logging parameter, metrik, dan artefak model
\end{itemize}

Validasi Statistik:

Pengujian hipotesis untuk memvalidasi signifikansi perbedaan antar konfigurasi (sebagaimana diuraikan Bab III):
\begin{itemize}[nosep]
    \item Uji: \textit{Independent samples t-test} (asumsi distribusi normal)
    \item Signifikansi: $\alpha = 0.05$ untuk hipotesis utama, $\alpha = 0.0167$ untuk \textit{multiple comparisons} (koreksi Bonferroni)
    \item Effect Size: Cohen's \textit{d} untuk mengukur \textit{magnitude} perbedaan
    \item Library: SciPy 1.11.4
    \item Kriteria: $p < \alpha$ DAN $d \geq 0.5$ (efek sedang hingga besar)
\end{itemize}

Semua detail implementasi ini memastikan transparansi metodologi dan reproduktifitas penuh sesuai standar riset \textit{open science}, dengan hasil evaluasi empiris dilaporkan secara komprehensif pada Bab V.

% \vspace{0.5em}
% \textit{Strategi pelatihan yang telah diuraikan pada subsection IV.4.1--IV.4.4 (curriculum learning, multi-component loss, discriminator mode, precision policy) membentuk metodologi inti framework GAN-HTR. Implementasi strategi ini memerlukan konfigurasi diskriminator dan sistem evaluasi yang tepat untuk memastikan adversarial equilibrium dan monitoring performa secara menyeluruh, yang akan diuraikan pada bagian berikutnya.}
% \vspace{0.5em}

% \textbf{Catatan:} \textit{Pipeline} pelatihan Recognizer yang digunakan sebagai komponen beku (metodologi, rasional, strategi curriculum learning) telah diuraikan pada Bab III Fase 2. Subsection IV.4 fokus pada strategi pelatihan framework GAN-HTR yang mengintegrasikan recognizer tersebut sebagai mekanisme \textit{guidance}.
% \newpage
% Print bibliography in APA style
% \printbibliography[heading=bibintoc]

